\chapter{Stable Brave New Motivic Homotopy Theory}
\label{chap:stable-brave-new-motivic}

\section{Purpose and standing stabilization conventions}
\label{sec:stable-purpose}

This chapter stabilizes the pointed brave new motivic category of the
preceding chapters.  For every quasi-compact quasi-separated spectral scheme
\(S\), Chapter~4 constructed a presentably symmetric monoidal pointed
motivic category
\[
    \HH_\bullet(S)
\]
under the smash product \(\wedge_S\), with tensor unit
\(\mathbb S^0_S\), and Chapter~5 supplied the open--closed localization
package, closed exchange, and homotopy purity.  In particular, for every
smooth closed pair
\[
    Z\lhook\joinrel\longrightarrow X
\]
over \(S\), there is a canonical equivalence
\[
    X/(X\setminus Z)
    \simeq
    \Th_S(\mathbf N_{Z/X})
\]
in \(\HH_\bullet(S)\), where
\[
    \mathbf N_{Z/X}
    =
    \mathbb V_Z(N^*_{Z/X})
\]
is the geometric normal bundle.  Equivalently, under the coordinate-module
convention for Thom spaces,
\[
    \Th_S(\mathbf N_{Z/X})
    =
    \Th_S(N^*_{Z/X}).
\]

The stabilization object is the Thom space of the trivial line bundle.  Write
\[
    T_S
    :=
    \Th_S(\mathcal O_S)
    \in
    \HH_\bullet(S).
\]
The stable brave new motivic category is the presentably symmetric monoidal periodization
\[
    \SH(S)
    :=
    \HH_\bullet(S)[T_S^{-1}].
\]
The principal purpose of the chapter is to construct this category together with its base
functoriality, smooth direct image, stable localization, and Thom twists.  The final output is a stable
$(\ast,\sharp,\otimes)$-formalism.

All conventions of the preceding chapters remain in force.  In particular, all categories are
$\infty$-categories, all coherence data are retained, and sheaves are not hypercompleted unless
explicitly stated otherwise.  Unless a statement explicitly concerns a more general base, every base
spectral scheme is assumed quasi-compact and quasi-separated.  Smooth morphisms are smooth in the
differential sense of Chapter~3 and are always assumed of finite presentation when smooth direct
image is invoked.

\begin{convention}[Scope of the stable coefficient system]
\label{conv:stable-coefficient-scope}
Whenever an expression of the form
\[
    \HH_\bullet(X),
    \qquad
    \SH(X),
    \qquad
    \Pic(\SH(X)),
\]
or a base-change, direct-image, or Thom-twist functor between such categories
occurs, the spectral scheme \(X\) is understood to be quasi-compact and
quasi-separated.

Statements involving only relative spectra, relative linear spaces,
cotangent complexes, zero sections, or finite locally free modules may be
stated for arbitrary spectral schemes.  In such statements, the
quasi-compact quasi-separated condition is imposed only when the associated
motivic coefficient category is invoked.
\end{convention}

\subsection{Pointed unstable input}

For a smooth spectral $S$-scheme $X$ of finite presentation, write
\[
    \Mot_{S,+}(X)
    \in
    \HH_\bullet(S)
\]
for its pointed motivic space.  The category $\HH_\bullet(S)$ is generated under colimits by these
objects.  For every morphism $f:T\to S$, there is a pointed adjunction
\[
    f^*:\HH_\bullet(S)
    \rightleftarrows
    \HH_\bullet(T):f_*,
\]
where $f^*$ is colimit-preserving and strong symmetric monoidal.  For every smooth morphism
$p:X\to S$ of finite presentation, there is a pointed adjunction
\[
    p_\sharp:\HH_\bullet(X)
    \rightleftarrows
    \HH_\bullet(S):p^*.
\]
These adjunctions satisfy pointed smooth base change and the pointed smooth projection formula.

For a complementary pair
\[
    Z\xrightarrow{i}S\xleftarrow{j}U,
\]
where $i$ is closed and $j$ is the quasi-compact open complement, Chapter~5 constructed
\[
    i^*:\HH_\bullet(S)
    \rightleftarrows
    \HH_\bullet(Z):i_*,
\]
\[
    j_\sharp:\HH_\bullet(U)
    \rightleftarrows
    \HH_\bullet(S):j^*
    \rightleftarrows
    \HH_\bullet(U):j_*,
\]
and a right adjoint
\[
    i_*:\HH_\bullet(Z)
    \rightleftarrows
    \HH_\bullet(S):i^!.
\]
The pointed localization sequences, closed base change, the closed projection formula, and
smooth--closed exchange will be stabilized below.

\subsection{The chapter-level output}

The constructions of the chapter culminate in the following data.

\begin{enumerate}
    \item A compactly generated stable presentably symmetric monoidal category $\SH(S)$.

    \item A strong symmetric monoidal stabilization functor
    \[
        \Sigma_T^\infty:\HH_\bullet(S)\longrightarrow\SH(S)
    \]
    with right adjoint $\Omega_T^\infty$.

    \item For every morphism $f:T\to S$, an adjunction
    \[
        f^*:\SH(S)\rightleftarrows\SH(T):f_*.
    \]

    \item For every smooth morphism $p:X\to S$ of finite presentation, an adjunction
    \[
        p_\sharp:\SH(X)\rightleftarrows\SH(S):p^*.
    \]

    \item For every virtual vector bundle $\xi\in K(S)$, a tensor-invertible Thom twist
    \[
        \Sigma^\xi:\SH(S)\xrightarrow{\simeq}\SH(S).
    \]

    \item Stable smooth base change, the smooth projection formula, stable localization,
    closed base change, the closed projection formula, and smooth--closed exchange.

    \item Stable homotopy purity.  For every smooth closed pair
    \[
        Z\lhook\joinrel\longrightarrow X
    \]
    over \(S\), with structural morphism \(z:Z\to S\), there is a canonical
    equivalence
    \[
        \Sigma_{T,S}^\infty
        \bigl(
          X/(X\setminus Z)
        \bigr)
        \simeq
        z_\sharp
        \operatorname{th}_Z
        \bigl(
          N^*_{Z/X}
        \bigr).
    \]

    \item Stable derived nilpotent invariance and comparison with ordinary stable motivic homotopy
    theory over the classical truncation.
\end{enumerate}

\section{Thom spaces as a symmetric monoidal construction}
\label{sec:stable-thom-spaces}

\subsection{Relative linear spaces and zero sections}

Fix a spectral scheme $X$ and a finite locally free quasi-coherent $\mathcal O_X$-module $E$.
Recall the relative linear space
\[
    \mathbb V_X(E)
    :=
    \operatorname{Spec}_X
    \operatorname{Sym}_{\mathcal O_X}(E).
\]
The augmentation
\[
    \operatorname{Sym}_{\mathcal O_X}(E)
    \longrightarrow
    \mathcal O_X
\]
induces the zero section
\[
    s_E:X\lhook\joinrel\longrightarrow\mathbb V_X(E).
\]
The structure morphism
\[
    \pi_E:\mathbb V_X(E)\longrightarrow X
\]
is smooth and of finite presentation.  The composite $\pi_Es_E$ is the identity of $X$.

\begin{proposition}[Base change for relative linear spaces]
\label{prop:linear-space-base-change}
Let $f:Y\to X$ be a morphism of spectral schemes.  There is a canonical Cartesian square
\[
\begin{tikzcd}[column sep=large,row sep=large]
    \mathbb V_Y(f^*E) \arrow[r] \arrow[d,"\pi_{f^*E}"']
        & \mathbb V_X(E) \arrow[d,"\pi_E"] \\
    Y \arrow[r,"f"'] & X,
\end{tikzcd}
\]
and the upper horizontal morphism carries the zero section of $f^*E$ to the pullback of the zero
section of $E$.
\end{proposition}

\begin{proof}
Relative spectrum commutes with base change, and the relative symmetric algebra satisfies
\[
    f^*\operatorname{Sym}_{\mathcal O_X}(E)
    \simeq
    \operatorname{Sym}_{\mathcal O_Y}(f^*E).
\]
The asserted Cartesian square follows from the base-change theorem for relative spectra.  The
augmentation of the pulled-back symmetric algebra is the pullback of the original augmentation, so
the assertion about zero sections follows as well.
\end{proof}

\begin{proposition}[Conormal module of the zero section]
\label{prop:zero-section-conormal}
The zero section
\[
    s_E:X\lhook\joinrel\longrightarrow\mathbb V_X(E)
\]
is a quasi-smooth closed immersion.  Its conormal module is canonically
equivalent to \(E\):
\[
    N^*_{X/\mathbb V_X(E)}
    \simeq
    E.
\]
Consequently, its normal module is canonically equivalent to \(E^\vee\):
\[
    N^{\mathrm{mod}}_{X/\mathbb V_X(E)}
    \simeq
    E^\vee.
\]
Under the convention
\[
    \mathbf N_{X/\mathbb V_X(E)}
    :=
    \mathbb V_X
    \bigl(
      (N^{\mathrm{mod}}_{X/\mathbb V_X(E)})^\vee
    \bigr),
\]
there is a canonical equivalence of geometric normal bundles
\[
    \mathbf N_{X/\mathbb V_X(E)}
    \simeq
    \mathbb V_X(E).
\]
\end{proposition}

\begin{proof}
The assertion is local on \(X\).  Write
\[
    X=\operatorname{Spec}(A)
\]
and let \(M\) be the finite locally free \(A\)-module corresponding to
\(E\).  Put
\[
    B:=\operatorname{Sym}_A(M).
\]
The zero section is represented by the augmentation
\[
    \varepsilon:B\longrightarrow A.
\]

We first verify that this morphism represents a closed immersion of finite
presentation.  After a Zariski localization on \(X\), choose an equivalence
\[
    M\simeq A^{\oplus r}.
\]
Then
\[
    B\simeq
    \operatorname{Sym}_A(A^{\oplus r}),
\]
and on degree-zero homotopy the augmentation is the ordinary polynomial
augmentation
\[
    \pi_0(A)[t_1,\ldots,t_r]
    \longrightarrow
    \pi_0(A),
    \qquad
    t_a\longmapsto0.
\]
It is surjective, so
\[
    \operatorname{Spec}(A)
    \longrightarrow
    \operatorname{Spec}(B)
\]
is a closed immersion.

The free-algebra calculation gives
\[
    L_{B/A}
    \simeq
    B\otimes_A M.
\]
Apply the transitivity triangle to the composite
\[
    A\longrightarrow B\longrightarrow A.
\]
Since the composite is the identity of \(A\), there is a cofibre sequence
\[
    A\otimes_B L_{B/A}
    \longrightarrow
    L_{A/A}
    \longrightarrow
    L_{A/B}.
\]
Using
\[
    L_{A/A}\simeq0,
\]
rotation gives
\[
    L_{A/B}
    \simeq
    \bigl(A\otimes_B L_{B/A}\bigr)[1].
\]
Therefore
\[
\begin{aligned}
    L_{A/B}
    &\simeq
    \bigl(
      A\otimes_B(B\otimes_A M)
    \bigr)[1] \\
    &\simeq
    M[1].
\end{aligned}
\]
In particular, \(L_{A/B}\) is perfect.  Moreover,
\(\pi_0(A)\) is finitely presented as a \(\pi_0(B)\)-algebra, since locally
it is the quotient
\[
    \pi_0(B)/(t_1,\ldots,t_r).
\]
The finite-presentation criterion detected by the cotangent complex therefore
shows that
\[
    B\longrightarrow A
\]
is of finite presentation.

Consequently, the zero section is a closed immersion of finite presentation.
Its conormal module is
\[
\begin{aligned}
    N^*_{X/\mathbb V_X(E)}
    &=
    L_{X/\mathbb V_X(E)}[-1] \\
    &\simeq
    E.
\end{aligned}
\]
Since \(E\) is finite locally free, the zero section is quasi-smooth.

The construction is natural under localization in \(A\), so the affine
equivalences glue canonically.  Dualizing gives
\[
    N^{\mathrm{mod}}_{X/\mathbb V_X(E)}
    \simeq
    E^\vee.
\]
Finite-projective biduality then gives
\[
\begin{aligned}
    \mathbf N_{X/\mathbb V_X(E)}
    &=
    \mathbb V_X
    \bigl(
      (N^{\mathrm{mod}}_{X/\mathbb V_X(E)})^\vee
    \bigr) \\
    &\simeq
    \mathbb V_X(E).
\end{aligned}
\]
\end{proof}

\subsection{Intrinsic and relative Thom spaces}

From this point onward in the Thom-space discussion, whenever an object is
placed in \(\HH_\bullet(X)\), \(\SH(X)\), or
\(\Pic(\SH(X))\), the spectral scheme \(X\) is assumed quasi-compact and
quasi-separated.  The preceding statements concerning relative spectra,
zero sections, and conormal modules remain valid without this hypothesis.

\begin{definition}[Intrinsic Thom space]
\label{def:intrinsic-thom}
Let \(X\) be a quasi-compact quasi-separated spectral scheme, and let \(E\)
be a finite locally free quasi-coherent \(\mathcal O_X\)-module.  The Thom
space of \(E\) over \(X\) is the pointed motivic quotient
\[
    \Th_X(E)
    :=
    \mathbb V_X(E)/
    \bigl(
      \mathbb V_X(E)\setminus X
    \bigr)
    \in
    \HH_\bullet(X),
\]
where \(X\) is embedded by the zero section.
\end{definition}

\begin{proposition}[Base change for Thom spaces]
\label{prop:thom-base-change-unstable}
Let $f:Y\to X$ be a morphism of quasi-compact quasi-separated spectral schemes.  For every finite
locally free $E$ on $X$, there is a canonical equivalence
\[
    f^*\Th_X(E)
    \simeq
    \Th_Y(f^*E)
\]
in $\HH_\bullet(Y)$.  These equivalences are coherent in $f$ and natural under equivalences of
vector bundles.
\end{proposition}

\begin{proof}
Pointed inverse image preserves colimits.  Proposition~\ref{prop:linear-space-base-change}
identifies the pullback of $\mathbb V_X(E)$ with $\mathbb V_Y(f^*E)$, and the pullback of the open
complement of the zero section with the corresponding open complement over $Y$.  Applying $f^*$
to the quotient defining $\Th_X(E)$ therefore gives the asserted equivalence.  Coherence follows
from coherent base change for relative spectra and pointed inverse image.
\end{proof}

Let
\[
    p:X\longrightarrow S
\]
be smooth and of finite presentation, and let \(E\) be finite locally free on
\(X\).  The structure morphism
\[
    \pi_E:\mathbb V_X(E)\longrightarrow X
\]
is smooth and of finite presentation, so the composite
\[
    \mathbb V_X(E)\longrightarrow S
\]
is smooth and of finite presentation.

By Proposition~\ref{prop:zero-section-conormal}, the zero section
\[
    s_E:X\hookrightarrow\mathbb V_X(E)
\]
is a quasi-smooth closed immersion.  In particular, it is a closed immersion
of finite presentation.  Its open complement
\[
    \mathbb V_X(E)\setminus X
    \longrightarrow
    \mathbb V_X(E)
\]
is therefore quasi-compact and open.  Since open immersions are smooth and
smooth morphisms are stable under composition, the complement
\[
    \mathbb V_X(E)\setminus X
\]
is also a smooth spectral \(S\)-scheme of finite presentation.  Hence both
terms in the following quotient define representable motivic spaces over
\(S\).

\begin{definition}[Relative Thom space]
\label{def:relative-thom-motive}
For \(p:X\to S\) smooth of finite presentation and \(E\) finite locally free
on \(X\), define
\[
    \Th_S(E)
    :=
    \Mot_S\bigl(\mathbb V_X(E)\bigr)
    /
    \Mot_S\bigl(\mathbb V_X(E)\setminus X\bigr)
    \in
    \HH_\bullet(S).
\]
\end{definition}

\begin{proposition}[Smooth direct image of intrinsic Thom spaces]
\label{prop:relative-thom-p-sharp}
Let $p:X\to S$ be smooth of finite presentation.  For every finite locally free module $E$ on $X$,
there is a canonical equivalence
\[
    p_\sharp\Th_X(E)
    \simeq
    \Th_S(E).
\]
\end{proposition}

\begin{proof}
Smooth direct image carries the pointed motive of a smooth $X$-scheme to the pointed motive of the
same spectral scheme viewed over $S$.  It preserves colimits.  Applying $p_\sharp$ to the quotient
that defines $\Th_X(E)$ therefore produces the quotient of the two corresponding smooth motives
over $S$.
\end{proof}

\subsection{Multiplicativity}

\begin{proposition}[Whitney sum formula for Thom spaces]
\label{prop:unstable-thom-whitney}
For finite locally free modules $E$ and $F$ on $X$, there is a canonical equivalence
\[
    \Th_X(E\oplus F)
    \simeq
    \Th_X(E)\wedge_X\Th_X(F).
\]
The equivalence is associative, symmetric, and unital with respect to direct sum.
\end{proposition}

\begin{proof}
There is a canonical equivalence of relative linear spaces
\[
    \mathbb V_X(E\oplus F)
    \simeq
    \mathbb V_X(E)\times_X\mathbb V_X(F).
\]
Under this equivalence, the complement of the joint zero section is the union
\[
\begin{aligned}
    &\bigl(\mathbb V_X(E)\setminus X\bigr)
        \times_X\mathbb V_X(F) \\
    &\qquad\cup
    \mathbb V_X(E)\times_X
        \bigl(\mathbb V_X(F)\setminus X\bigr).
\end{aligned}
\]
The smash product of the two pointed quotients is the quotient of
$\mathbb V_X(E)\times_X\mathbb V_X(F)$ by precisely this union.  Hence the pushout formula for the
smash product gives the asserted equivalence.

The construction is induced by the canonical equivalence of symmetric algebras
\[
    \operatorname{Sym}(E\oplus F)
    \simeq
    \operatorname{Sym}(E)\otimes\operatorname{Sym}(F).
\]
Its associativity, symmetry, and unit coherences are therefore inherited from direct sum and the
symmetric monoidal relative-spectrum construction.
\end{proof}

\begin{corollary}[The unstable Thom functor]
\label{cor:unstable-thom-functor}
The assignment $E\mapsto\Th_X(E)$ refines to a symmetric monoidal functor
\[
    \Th_X:
    \bigl(\Vect(X)^\simeq,\oplus\bigr)
    \longrightarrow
    \bigl(\HH_\bullet(X),\wedge_X\bigr).
\]
It is natural under pullback in $X$.
\end{corollary}

\begin{proof}
Proposition~\ref{prop:unstable-thom-whitney} supplies the binary monoidal comparison.  The zero
bundle has relative linear space $X$ and empty zero-section complement, so its Thom space is
$\mathbb S^0_X$.  The coherence assertions follow from the final paragraph of the proof of
Proposition~\ref{prop:unstable-thom-whitney}.  Naturality is
Proposition~\ref{prop:thom-base-change-unstable}.
\end{proof}

\begin{definition}[The motivic Thom sphere]
\label{def:T-sphere}
Define
\[
    T_X
    :=
    \Th_X(\mathcal O_X).
\]
Equivalently,
\[
    T_X
    \simeq
    \mathbb A^1_X/\bigl(\mathbb A^1_X\setminus0_X\bigr).
\]
\end{definition}

\begin{corollary}[Trivial bundles]
\label{cor:trivial-thom-power}
For every integer $n\geq0$, there is a canonical equivalence
\[
    \Th_X(\mathcal O_X^{\oplus n})
    \simeq
    T_X^{\wedge n}.
\]
\end{corollary}

\begin{proof}
Iterate Proposition~\ref{prop:unstable-thom-whitney}.
\end{proof}

\section{The motivic circle decomposition of \texorpdfstring{$T$}{T}}
\label{sec:T-circle-decomposition}

\subsection{The punctured brave new affine line}

Let
\[
    \Gm{}_X
    :=
    \mathbb A^1_X\setminus0_X.
\]
The unit section $1_X:X\to\Gm{}_X$ makes $\Gm{}_X$ a pointed motivic space.  Base change of the
brave new affine line and of principal open complements gives
\[
    f^*\Gm{}_X
    \simeq
    \Gm{}_Y
\]
for every $f:Y\to X$.

The brave new affine line is contractible in $\HH(X)$ by construction of the motivic localization:
its structural projection induces an equivalence
\[
    \Mot_X(\mathbb A^1_X)
    \simeq
    \Mot_X(X).
\]
After adjoining disjoint basepoints, this becomes
\[
    \Mot_{X,+}(\mathbb A^1_X)
    \simeq
    \mathbb S^0_X.
\]

\begin{proposition}[Circle decomposition of the Thom sphere]
\label{prop:T-circle-Gm}
There is a canonical equivalence
\[
    T_X
    \simeq
    S^1\wedge\Gm{}_X
\]
in \(\HH_\bullet(X)\), where \(S^1\) denotes the simplicial suspension
circle and \(\Gm{}_X\) is pointed by the unit section
\[
    1_X:X\longrightarrow\Gm{}_X.
\]
The equivalence is compatible with base change in \(X\).
\end{proposition}

\begin{proof}
Point \(\mathbb A^1_X\) by the unit section
\[
    1_X:X\longrightarrow\mathbb A^1_X.
\]
The inclusion of the punctured affine line is then a morphism of pointed
motivic spaces
\[
    (\Gm{}_X,1_X)
    \longrightarrow
    (\mathbb A^1_X,1_X).
\]
Its pointed cofiber is canonically the quotient
\[
    \mathbb A^1_X/
    \bigl(\mathbb A^1_X\setminus0_X\bigr),
\]
because the basepoint \(1_X\) lies in
\(\mathbb A^1_X\setminus0_X\).  Hence there is a canonical equivalence
\[
    T_X
    \simeq
    \operatorname{cofib}
    \bigl(
      (\Gm{}_X,1_X)
      \longrightarrow
      (\mathbb A^1_X,1_X)
    \bigr).
\]

The structural projection
\[
    \pi:\mathbb A^1_X\longrightarrow X
\]
is inverted by motivic affine-line localization.  Since
\[
    \pi 1_X=\id_X,
\]
the unit section is an inverse equivalence to \(\pi\) in \(\HH(X)\).
With the indicated basepoint, this gives an equivalence
\[
    (\mathbb A^1_X,1_X)
    \simeq
    0_X
\]
in \(\HH_\bullet(X)\), where \(0_X\) denotes the zero object.

It follows that
\[
\begin{aligned}
    T_X
    &\simeq
    \operatorname{cofib}
    \bigl(
      (\Gm{}_X,1_X)\longrightarrow0_X
    \bigr) \\
    &\simeq
    \Sigma(\Gm{}_X,1_X) \\
    &\simeq
    S^1\wedge\Gm{}_X.
\end{aligned}
\]

The zero section, unit section, punctured affine line, and affine-line
contraction are all preserved by base change.  The resulting equivalence is
therefore compatible with base change in \(X\).
\end{proof}

\begin{lemma}[Factors of an invertible tensor product]
\label{lem:factors-invertible}
Let $\mathcal C$ be a symmetric monoidal $\infty$-category.  If $A\otimes B$ is tensor-invertible,
then both $A$ and $B$ are tensor-invertible.
\end{lemma}

\begin{proof}
Let $C$ be an inverse of $A\otimes B$.  Then $B\otimes C$ is an inverse of $A$, since
\[
    A\otimes(B\otimes C)
    \simeq
    (A\otimes B)\otimes C
    \simeq
    \mathbf 1.
\]
Symmetry gives the corresponding equivalence in the opposite order.  The same argument, with
$A$ and $B$ interchanged, gives an inverse of $B$.
\end{proof}

\begin{corollary}[Consequences of \(T\)-invertibility]
\label{cor:T-inversion-circle}
Let
\[
    F:
    \HH_\bullet(X)
    \longrightarrow
    \mathcal D
\]
be a strong symmetric monoidal functor to a symmetric monoidal
\(\infty\)-category \(\mathcal D\).  If \(F(T_X)\) is tensor-invertible,
then both
\[
    F(S^1_X)
    \qquad\text{and}\qquad
    F(\Gm{}_X)
\]
are tensor-invertible in \(\mathcal D\).
\end{corollary}

\begin{proof}
Proposition~\ref{prop:T-circle-Gm} gives
\[
    T_X
    \simeq
    S^1_X\wedge_X\Gm{}_X.
\]
Since \(F\) is strong symmetric monoidal,
\[
    F(T_X)
    \simeq
    F(S^1_X)\otimes F(\Gm{}_X).
\]
Apply Lemma~\ref{lem:factors-invertible}.
\end{proof}

\begin{definition}[Symmetric tensor object]
\label{def:symmetric-tensor-object}
Let $\mathcal C$ be a symmetric monoidal $\infty$-category.  An object $X\in\mathcal C$ is
symmetric if the cyclic permutation
\[
    (123):X^{\otimes3}\longrightarrow X^{\otimes3}
\]
is homotopic to the identity.
\end{definition}

\begin{proposition}[Symmetry of the motivic Thom sphere]
\label{prop:T-symmetric}
For every quasi-compact quasi-separated spectral scheme \(X\), the object
\[
    T_X\in\HH_\bullet(X)
\]
is symmetric.
\end{proposition}

\begin{proof}
The Whitney-sum equivalence identifies
\[
    T_X^{\wedge3}
    \simeq
    \Th_X(\mathcal O_X^{\oplus3}).
\]
Under this identification, the cyclic permutation of the three tensor
factors is induced by the cyclic permutation matrix
\[
    P=
    \begin{pmatrix}
        0&1&0\\
        0&0&1\\
        1&0&0
    \end{pmatrix}.
\]

For \(i\neq j\), let \(E_{ij}(a)\) denote the elementary matrix with
off-diagonal entry \(a\) in position \((i,j)\).  There is an identity
\[
    P
    =
    E_{23}(1)E_{32}(-1)E_{23}(1)
    E_{12}(1)E_{21}(-1)E_{12}(1).
\]
Let \(t\) be the coordinate on the brave new affine line and define
\[
    H(t)
    :=
    E_{23}(t)E_{32}(-t)E_{23}(t)
    E_{12}(t)E_{21}(-t)E_{12}(t).
\]
Every elementary factor is invertible, with inverse obtained by replacing
\(t\) by \(-t\).  Hence \(H(t)\) is an automorphism of the trivial
rank-three bundle over
\[
    X\times\mathbb A^1.
\]
It satisfies
\[
    H(0)=I,
    \qquad
    H(1)=P.
\]

The corresponding automorphism of the relative linear space preserves the
zero section and therefore also its open complement.  Passing to the pointed
quotient gives an \(\mathbb A^1\)-homotopy between the identity of
\[
    \Th_X(\mathcal O_X^{\oplus3})
\]
and the cyclic permutation.  Affine-line localization identifies the two
endpoint maps.  Therefore the cyclic permutation of \(T_X^{\wedge3}\) is
homotopic to the identity.
\end{proof}

\section{Symmetric monoidal periodization}
\label{sec:abstract-periodization}

This section records the formal inversion results used throughout the chapter.  They are applied in
presentably symmetric monoidal pointed categories, and all functors are required to preserve small
colimits.  Write $\Pr^L_*$ for the symmetric monoidal $\infty$-category of presentable pointed
$\infty$-categories and basepoint-preserving colimit-preserving functors.

\subsection{Formal inversion of an object}

\begin{lemma}[Free pointed presentably symmetric monoidal completion]
\label{lem:free-pointed-monoidal-completion}
Let
\[
    \mathcal A^\otimes
\]
be a small symmetric monoidal \(\infty\)-category.  Put
\[
    \mathcal P_*(\mathcal A)
    :=
    \Fun(\mathcal A^{\op},\mathcal S_*).
\]
Equip this category with pointed Day convolution.  Then
\(\mathcal P_*(\mathcal A)\) is a presentably symmetric monoidal pointed
\(\infty\)-category whose tensor product preserves small colimits separately
in each variable.

The pointed Yoneda functor
\[
    y_+:
    \mathcal A
    \longrightarrow
    \mathcal P_*(\mathcal A),
    \qquad
    a\longmapsto
    \Map_{\mathcal A}(-,a)_+,
\]
is strong symmetric monoidal.

For every presentably symmetric monoidal pointed \(\infty\)-category
\(\mathcal D\), restriction along \(y_+\) induces an equivalence
\[
    \Fun^{L,\otimes}
    \bigl(
      \mathcal P_*(\mathcal A),
      \mathcal D
    \bigr)
    \xrightarrow{\simeq}
    \Fun^\otimes(\mathcal A,\mathcal D).
\]
The equivalence is natural in both \(\mathcal A\) and \(\mathcal D\).
\end{lemma}

\begin{proof}
There is a canonical equivalence
\[
    \Fun(\mathcal A^{\op},\mathcal S_*)
    \simeq
    \Fun(\mathcal A^{\op},\mathcal S)_*
\]
between pointed-space-valued presheaves and pointed objects in the ordinary
presheaf category.

The symmetric monoidal structure on \(\mathcal A\) and the smash product on
\(\mathcal S_*\) determine the pointed Day-convolution product on
\(\mathcal P_*(\mathcal A)\).  It is characterized by the following two
properties:
\begin{enumerate}
    \item the pointed Yoneda functor \(y_+\) is strong symmetric monoidal;
    \item the tensor product preserves small colimits separately in each
    variable.
\end{enumerate}
Since colimits of pointed presheaves are computed objectwise and the smash
product on pointed spaces preserves colimits separately, pointed Day
convolution also preserves colimits separately.

Let \(\mathcal D\) be a presentably symmetric monoidal pointed
\(\infty\)-category.  Given a symmetric monoidal functor
\[
    F:
    \mathcal A
    \longrightarrow
    \mathcal D,
\]
its left Kan extension along \(y_+\) exists and preserves small colimits:
\[
    \widetilde F:
    \mathcal P_*(\mathcal A)
    \longrightarrow
    \mathcal D.
\]
The Day-convolution universal property equips \(\widetilde F\) with a
canonical strong symmetric monoidal structure.  By construction,
\[
    \widetilde F y_+
    \simeq
    F.
\]

Conversely, every colimit-preserving strong symmetric monoidal functor
\[
    G:
    \mathcal P_*(\mathcal A)
    \longrightarrow
    \mathcal D
\]
is determined by its restriction to pointed representables, because the
objects \(y_+(a)\) generate \(\mathcal P_*(\mathcal A)\) under small
colimits.  Its monoidal structure is likewise determined by the monoidal
structure on those generators, since both tensor products preserve colimits
separately.

Thus left Kan extension and restriction are inverse equivalences:
\[
    \Fun^{L,\otimes}
    \bigl(
      \mathcal P_*(\mathcal A),
      \mathcal D
    \bigr)
    \simeq
    \Fun^\otimes(\mathcal A,\mathcal D).
\]
Naturality follows from the naturality of pointed Yoneda, Day convolution,
and left Kan extension.
\end{proof}

\begin{theorem}[Presentable symmetric monoidal periodization]
\label{thm:presentable-periodization}
Let \(\mathcal C\) be a presentably symmetric monoidal pointed
\(\infty\)-category and let \(X\in\mathcal C\).  There exist a presentably
symmetric monoidal pointed \(\infty\)-category
\[
    \mathcal C[X^{-1}]
\]
and a colimit-preserving strong symmetric monoidal functor
\[
    \operatorname{per}_X:
    \mathcal C
    \longrightarrow
    \mathcal C[X^{-1}]
\]
such that:
\begin{enumerate}
    \item the object
    \[
        \operatorname{per}_X(X)
    \]
    is tensor-invertible;

    \item for every presentably symmetric monoidal pointed
    \(\infty\)-category \(\mathcal D\), restriction along
    \(\operatorname{per}_X\) induces an equivalence
    \[
    \begin{aligned}
        \Fun^{L,\otimes}
        \bigl(
          \mathcal C[X^{-1}],
          \mathcal D
        \bigr)
        \simeq
        \bigl\{
          F\in
          \Fun^{L,\otimes}(\mathcal C,\mathcal D)
          \mid
          F(X)\text{ is tensor-invertible}
        \bigr\};
    \end{aligned}
    \]

    \item the pair
    \[
        \bigl(
          \mathcal C[X^{-1}],
          \operatorname{per}_X
        \bigr)
    \]
    is unique up to a contractible space of equivalences compatible with
    the functor from \(\mathcal C\).
\end{enumerate}
\end{theorem}

\begin{proof}
Let
\[
    \mathbb F^\otimes
\]
be the free small symmetric monoidal \(\infty\)-category on one object
\(x\).  Let
\[
    \mathbb F^\otimes[x^{-1}]
\]
be its symmetric monoidal localization obtained by adjoining an inverse to
\(x\).  Thus, for every symmetric monoidal \(\infty\)-category
\(\mathcal E\), restriction induces an equivalence
\[
\begin{aligned}
    \Fun^\otimes
    \bigl(
      \mathbb F[x^{-1}],
      \mathcal E
    \bigr)
    \simeq
    \bigl\{
      F\in\Fun^\otimes(\mathbb F,\mathcal E)
      \mid
      F(x)\text{ is tensor-invertible}
    \bigr\}.
\end{aligned}
\]

Put
\[
    \mathcal F
    :=
    \mathcal P_*(\mathbb F),
    \qquad
    \mathcal F^{\mathrm{inv}}
    :=
    \mathcal P_*
    \bigl(
      \mathbb F[x^{-1}]
    \bigr).
\]
The localization
\[
    \mathbb F
    \longrightarrow
    \mathbb F[x^{-1}]
\]
induces a colimit-preserving strong symmetric monoidal functor
\[
    \lambda:
    \mathcal F
    \longrightarrow
    \mathcal F^{\mathrm{inv}}.
\]

The chosen object \(X\in\mathcal C\) determines a symmetric monoidal
functor
\[
    \mathbb F
    \longrightarrow
    \mathcal C,
    \qquad
    x\longmapsto X.
\]
By Lemma~\ref{lem:free-pointed-monoidal-completion}, this extends uniquely,
up to a contractible space of choices, to a colimit-preserving strong
symmetric monoidal functor
\[
    u_X:
    \mathcal F
    \longrightarrow
    \mathcal C.
\]

Define
\[
    \mathcal C[X^{-1}]
    :=
    \mathcal C
    \otimes_{\mathcal F}
    \mathcal F^{\mathrm{inv}}
\]
in
\[
    \CAlg(\Pr^L_*).
\]
Let
\[
    \operatorname{per}_X:
    \mathcal C
    \longrightarrow
    \mathcal C[X^{-1}]
\]
be the canonical morphism.

The image of \(X\) in the relative tensor product is identified with the
image of \(x\) from \(\mathcal F^{\mathrm{inv}}\).  Since \(x\) is
tensor-invertible there,
\[
    \operatorname{per}_X(X)
\]
is tensor-invertible.

Let \(\mathcal D\) be a presentably symmetric monoidal pointed
\(\infty\)-category.  Mapping out of the relative tensor product gives a
pullback of mapping spaces
\[
\begin{aligned}
    \Fun^{L,\otimes}
    \bigl(
      \mathcal C[X^{-1}],
      \mathcal D
    \bigr)
    \simeq
    \Fun^{L,\otimes}(\mathcal C,\mathcal D)
    \times_{\Fun^{L,\otimes}(\mathcal F,\mathcal D)}
    \Fun^{L,\otimes}
    \bigl(
      \mathcal F^{\mathrm{inv}},
      \mathcal D
    \bigr).
\end{aligned}
\]

By Lemma~\ref{lem:free-pointed-monoidal-completion}, a functor
\[
    \mathcal F\longrightarrow\mathcal D
\]
is determined by the image of \(x\).  The localization property of
\(\mathbb F[x^{-1}]\) shows that such a functor extends across
\(\lambda\) exactly when that image is tensor-invertible.  When an
extension exists, its space of choices is contractible.

Consequently, the preceding pullback is canonically equivalent to
\[
    \bigl\{
      F\in
      \Fun^{L,\otimes}(\mathcal C,\mathcal D)
      \mid
      F(X)\text{ is tensor-invertible}
    \bigr\}.
\]
This proves the universal property.

Finally, any two pairs satisfying this universal property represent the
same functor on presentably symmetric monoidal pointed targets.  The
\(\infty\)-categorical Yoneda lemma therefore gives a contractible space of
equivalences between them, compatible with their functors from
\(\mathcal C\).
\end{proof}

\begin{lemma}[Symmetric-object telescope comparison]
\label{lem:symmetric-object-telescope-comparison}
Let \(\mathcal C\) be a presentably symmetric monoidal pointed
\(\infty\)-category, and let \(X\in\mathcal C\) be symmetric in the sense
that the cyclic permutation
\[
    (123):
    X^{\otimes3}
    \longrightarrow
    X^{\otimes3}
\]
is homotopic to the identity.

Put
\[
    L_X
    :=
    \operatorname{per}_X(X)
    \in
    \Pic\bigl(\mathcal C[X^{-1}]\bigr).
\]
For \(n\geq0\), write \(L_X^{-n}\) for the negative tensor power in the
Picard \(\infty\)-groupoid.  The space of choices involved in forming these
inverse powers is contractible.

Then there is a canonical equivalence in \(\Pr^L_*\)
\[
\begin{aligned}
    \operatorname*{colim}
    \left(
      \mathcal C
      \xrightarrow{X\otimes-}
      \mathcal C
      \xrightarrow{X\otimes-}
      \mathcal C
      \longrightarrow\cdots
    \right)
    \xrightarrow{\simeq}
    \mathcal C[X^{-1}].
\end{aligned}
\]
On the \(n\)-th stage, the comparison is canonically equivalent to
\[
    C
    \longmapsto
    L_X^{-n}
    \otimes
    \operatorname{per}_X(C).
\]
\end{lemma}

\begin{proof}
The formal inversion theorem for a symmetric object identifies the
underlying presentable category of
\[
    \mathcal C[X^{-1}]
\]
with the stabilization of the \(\mathcal C\)-module \(\mathcal C\) under
the action of \(X\).

When the cyclic permutation on \(X^{\otimes3}\) is homotopic to the
identity, the action of \(X\) on the sequential stabilization is an
equivalence.  The stabilization is therefore computed by the telescope
\[
    \mathcal C
    \xrightarrow{X\otimes-}
    \mathcal C
    \xrightarrow{X\otimes-}
    \mathcal C
    \longrightarrow\cdots.
\]
The universal property of formal inversion identifies this telescope with
the underlying presentable category of
\(\mathcal C[X^{-1}]\); see
\cite[Proposition~2.19 and Corollary~2.22]
{RobaloNoncommutativeMotivesI}.

Under this equivalence, an object \(C\) placed in the \(n\)-th stage is
transported to the periodized object
\(\operatorname{per}_X(C)\), together with \(n\) inverse applications of
the invertible action of \(X\).  Its image is therefore
\[
    L_X^{-n}
    \otimes
    \operatorname{per}_X(C).
\]
The uniqueness statement in the formal inversion theorem makes these
stagewise identifications canonical through a contractible space of
choices.
\end{proof}

\begin{theorem}[Symmetric telescope presentation]
\label{thm:periodization-telescope}
Let \(\mathcal C\) be a presentably symmetric monoidal pointed
\(\infty\)-category, and let \(X\in\mathcal C\) be symmetric.  Define
\[
    \operatorname{Tel}_X(\mathcal C)
    :=
    \operatorname*{colim}
    \left(
      \mathcal C
      \xrightarrow{X\otimes-}
      \mathcal C
      \xrightarrow{X\otimes-}
      \mathcal C
      \xrightarrow{X\otimes-}
      \cdots
    \right)
\]
in \(\Pr^L_*\).

There is a canonical equivalence
\[
    \Phi_X:
    \operatorname{Tel}_X(\mathcal C)
    \xrightarrow{\simeq}
    \mathcal C[X^{-1}]
\]
in \(\Pr^L_*\).  Equip
\(\operatorname{Tel}_X(\mathcal C)\) with the presentably symmetric
monoidal structure transported from \(\mathcal C[X^{-1}]\) across
\(\Phi_X\).

For this transported structure, the canonical functor
\[
    \iota_0:
    \mathcal C
    \longrightarrow
    \operatorname{Tel}_X(\mathcal C)
\]
is colimit-preserving and strong symmetric monoidal, and
\[
    \iota_0(X)
\]
is tensor-invertible.

If
\[
    \iota_n:
    \mathcal C
    \longrightarrow
    \operatorname{Tel}_X(\mathcal C)
\]
is the canonical functor from the \(n\)-th stage, then there are canonical
equivalences, relative to the transported monoidal structure,
\[
    \iota_m(C)\otimes\iota_n(D)
    \simeq
    \iota_{m+n}(C\otimes D).
\]

For every presentably symmetric monoidal pointed \(\infty\)-category
\(\mathcal D\), restriction along \(\iota_0\) induces an equivalence
\[
\begin{aligned}
    \Fun^{L,\otimes}
    \bigl(
      \operatorname{Tel}_X(\mathcal C),
      \mathcal D
    \bigr)
    \simeq
    \bigl\{
      F\in
      \Fun^{L,\otimes}(\mathcal C,\mathcal D)
      \mid
      F(X)\text{ is tensor-invertible}
    \bigr\}.
\end{aligned}
\]
\end{theorem}

\begin{proof}
Lemma~\ref{lem:symmetric-object-telescope-comparison} gives the canonical
equivalence
\[
    \Phi_X:
    \operatorname{Tel}_X(\mathcal C)
    \xrightarrow{\simeq}
    \mathcal C[X^{-1}]
\]
of presentable pointed \(\infty\)-categories.  Transport the presentably
symmetric monoidal structure of \(\mathcal C[X^{-1}]\) across
\(\Phi_X\).

Under this equivalence,
\[
    \Phi_X\iota_0
    \simeq
    \operatorname{per}_X.
\]
Therefore \(\iota_0\) is colimit-preserving and strong symmetric monoidal,
and \(\iota_0(X)\) is tensor-invertible.

Put
\[
    L_X=\operatorname{per}_X(X).
\]
Lemma~\ref{lem:symmetric-object-telescope-comparison} identifies
\[
    \Phi_X\bigl(\iota_m(C)\bigr)
    \simeq
    L_X^{-m}\otimes\operatorname{per}_X(C)
\]
and
\[
    \Phi_X\bigl(\iota_n(D)\bigr)
    \simeq
    L_X^{-n}\otimes\operatorname{per}_X(D).
\]
Consequently,
\[
\begin{aligned}
    \Phi_X
    \bigl(
      \iota_m(C)\otimes\iota_n(D)
    \bigr)
    &\simeq
    L_X^{-(m+n)}
    \otimes
    \operatorname{per}_X(C\otimes D) \\
    &\simeq
    \Phi_X
    \bigl(
      \iota_{m+n}(C\otimes D)
    \bigr).
\end{aligned}
\]
Since \(\Phi_X\) is an equivalence, this gives the displayed stage formula.
Its associativity, symmetry, unit, and higher coherence are inherited from
the transported symmetric monoidal structure.

Finally,
\[
\begin{aligned}
    \Fun^{L,\otimes}
    \bigl(
      \operatorname{Tel}_X(\mathcal C),
      \mathcal D
    \bigr)
    &\simeq
    \Fun^{L,\otimes}
    \bigl(
      \mathcal C[X^{-1}],
      \mathcal D
    \bigr).
\end{aligned}
\]
Apply Theorem~\ref{thm:presentable-periodization} to identify the right-hand
side with the required space of functors.
\end{proof}

\subsection{Linear functors and periodization}

\begin{definition}[Linear functor relative to two objects]
\label{def:XY-linear}
Let $(\mathcal C,X)$ and $(\mathcal D,Y)$ be presentably symmetric monoidal categories with chosen
objects.  A colimit-preserving functor $L:\mathcal C\to\mathcal D$ is $(X,Y)$-linear if it is equipped
with a coherent natural equivalence
\[
    L(X\otimes C)
    \simeq
    Y\otimes L(C).
\]
\end{definition}

\begin{proposition}[Periodization of linear functors]
\label{prop:linear-functor-periodization}
Suppose that $X$ and $Y$ are symmetric, and let $L:\mathcal C\to\mathcal D$ be
$(X,Y)$-linear.  Then there is a canonically induced
colimit-preserving functor
\[
    L^{\mathrm{per}}:
    \mathcal C[X^{-1}]
    \longrightarrow
    \mathcal D[Y^{-1}]
\]
and a canonical equivalence
\[
    L^{\mathrm{per}}\operatorname{per}_X
    \simeq
    \operatorname{per}_Y L.
\]
The construction is compatible with composition and with coherent natural transformations of
linear functors.
\end{proposition}

\begin{proof}
Using the sequential presentation of Theorem~\ref{thm:periodization-telescope}, the linearity
equivalence gives a morphism between the two sequential diagrams defining the periodizations.
Taking the colimit in $\Pr^L$ gives $L^{\mathrm{per}}$ and the displayed compatibility with the
periodization functors.  Composition of linear functors and coherent natural transformations define
the corresponding compositions and transformations of telescope diagrams, so they descend through
the same colimit construction.
\end{proof}

\begin{proposition}[Periodization of multilinear transformations]
\label{prop:multilinear-periodization}
For \(1\leq a\leq r\), let
\[
    (\mathcal C_a,X_a)
\]
be presentably symmetric monoidal categories equipped with symmetric objects,
and let
\[
    (\mathcal D,Y)
\]
be another such pair.

Suppose
\[
    F,G:
    \mathcal C_1\times\cdots\times\mathcal C_r
    \longrightarrow
    \mathcal D
\]
preserve small colimits separately in every variable.  Suppose further that,
for every \(a\), they are equipped with coherent linearity equivalences
\[
\begin{aligned}
    F(C_1,\ldots,X_a\otimes C_a,\ldots,C_r)
    &\simeq
    Y\otimes
    F(C_1,\ldots,C_a,\ldots,C_r), \\
    G(C_1,\ldots,X_a\otimes C_a,\ldots,C_r)
    &\simeq
    Y\otimes
    G(C_1,\ldots,C_a,\ldots,C_r),
\end{aligned}
\]
and that the linearity structures for distinct variables commute coherently.

Then \(F\) and \(G\) induce separately colimit-preserving functors
\[
    F^{\mathrm{per}},G^{\mathrm{per}}:
    \mathcal C_1[X_1^{-1}]
    \times\cdots\times
    \mathcal C_r[X_r^{-1}]
    \longrightarrow
    \mathcal D[Y^{-1}].
\]

Every natural transformation
\[
    \alpha:F\longrightarrow G
\]
compatible with the multilinearity structures induces a natural
transformation
\[
    \alpha^{\mathrm{per}}:
    F^{\mathrm{per}}
    \longrightarrow
    G^{\mathrm{per}}.
\]
If \(\alpha\) is an equivalence, then
\(\alpha^{\mathrm{per}}\) is an equivalence.

The construction is compatible with composition, and different orders of
periodizing the source variables are canonically coherently equivalent.
\end{proposition}

\begin{proof}
Write
\[
    \operatorname{per}_Y:
    \mathcal D
    \longrightarrow
    \mathcal D[Y^{-1}]
\]
for the periodization functor.  The object
\[
    \operatorname{per}_Y(Y)
\]
is tensor-invertible.  Consequently, periodizing
\(\mathcal D[Y^{-1}]\) once more at
\(\operatorname{per}_Y(Y)\) gives a canonical equivalence
\[
    \mathcal D[Y^{-1}]
    \bigl[
      \operatorname{per}_Y(Y)^{-1}
    \bigr]
    \simeq
    \mathcal D[Y^{-1}].
\]

We construct \(F^{\mathrm{per}}\) by successively periodizing the source
variables.  Begin with the composite
\[
    F^{(0)}
    :=
    \operatorname{per}_Y F:
    \mathcal C_1\times\cdots\times\mathcal C_r
    \longrightarrow
    \mathcal D[Y^{-1}].
\]
For fixed objects in the variables \(2,\ldots,r\), the functor
\[
    C_1
    \longmapsto
    F^{(0)}(C_1,C_2,\ldots,C_r)
\]
is
\[
    \bigl(
      X_1,\operatorname{per}_Y(Y)
    \bigr)\text{-linear}.
\]
Proposition~\ref{prop:linear-functor-periodization}, followed by the preceding
idempotence equivalence, therefore gives a functor
\[
    F^{(1)}:
    \mathcal C_1[X_1^{-1}]
    \times
    \mathcal C_2
    \times\cdots\times
    \mathcal C_r
    \longrightarrow
    \mathcal D[Y^{-1}].
\]

The construction in
Proposition~\ref{prop:linear-functor-periodization} is functorial in all
parameters and in coherent natural transformations.  Hence \(F^{(1)}\)
preserves colimits separately in every variable.  Moreover, the commuting
linearity data for \(F\) imply that, for every \(a\geq2\), the functor
\(F^{(1)}\) remains
\[
    \bigl(
      X_a,\operatorname{per}_Y(Y)
    \bigr)\text{-linear}
\]
in the \(a\)-th variable.

Periodizing the second variable gives
\[
    F^{(2)}:
    \mathcal C_1[X_1^{-1}]
    \times
    \mathcal C_2[X_2^{-1}]
    \times
    \mathcal C_3
    \times\cdots\times
    \mathcal C_r
    \longrightarrow
    \mathcal D[Y^{-1}].
\]
Continuing inductively gives
\[
    F^{(r)}:
    \mathcal C_1[X_1^{-1}]
    \times\cdots\times
    \mathcal C_r[X_r^{-1}]
    \longrightarrow
    \mathcal D[Y^{-1}].
\]
Define
\[
    F^{\mathrm{per}}
    :=
    F^{(r)}.
\]
The same construction gives \(G^{\mathrm{per}}\).

Now let
\[
    \alpha:F\longrightarrow G
\]
be compatible with all multilinearity structures.  After composing with
\(\operatorname{per}_Y\), Proposition~\ref{prop:linear-functor-periodization}
periodizes \(\alpha\) in the first variable.  Compatibility with the remaining
linearity structures allows the resulting transformation to be periodized in
the second variable, and so on.  After \(r\) steps one obtains
\[
    \alpha^{\mathrm{per}}:
    F^{\mathrm{per}}
    \longrightarrow
    G^{\mathrm{per}}.
\]
If \(\alpha\) is an equivalence, its inverse is compatible with the same
multilinearity data.  Periodizing the inverse and the two inverse homotopies
shows that \(\alpha^{\mathrm{per}}\) is an equivalence.

It remains to compare different orders of periodization.  For any ordering
of the variables, the resulting functor is characterized by the following
properties:

\begin{enumerate}
    \item its restriction along the product of the periodization functors is
    canonically equivalent to \(\operatorname{per}_YF\);

    \item it is separately colimit-preserving;

    \item it is compatible with the specified inverse actions of the
    \(X_a\) through the invertible object \(\operatorname{per}_Y(Y)\).
\end{enumerate}

Repeated use of the universal property of single-variable periodization shows
that the space of functors with these properties is contractible.  Therefore
different orders of periodization are canonically equivalent.  The same
contractible uniqueness gives the required higher coherence and compatibility
with composition.
\end{proof}

\begin{proposition}[Descent of adjunctions]
\label{prop:periodized-adjunction}
Suppose that $X$ and $Y$ are symmetric and that
\[
    L:\mathcal C\rightleftarrows\mathcal D:R
\]
is an adjunction.  Suppose that both $L$ and $R$ preserve colimits, that $L$ is $(X,Y)$-linear, that
$R$ is $(Y,X)$-linear, and that the two linearity equivalences are adjoint mates.  Then the
periodizations form an adjunction
\[
    L^{\mathrm{per}}:
    \mathcal C[X^{-1}]
    \rightleftarrows
    \mathcal D[Y^{-1}]:R^{\mathrm{per}}.
\]
If a square of such adjunctions and its Beck--Chevalley transformation are compatible with the
chosen linearity data, then the square, the transformation, and its adjoint mate descend to the
periodizations.  Equivalences descend to equivalences.
\end{proposition}

\begin{proof}
Proposition~\ref{prop:linear-functor-periodization} separately constructs $L^{\mathrm{per}}$ and
$R^{\mathrm{per}}$.  Compatibility of the linearity structures with adjunction makes the unit
\[
    \id_{\mathcal C}\longrightarrow RL
\]
and counit
\[
    LR\longrightarrow\id_{\mathcal D}
\]
coherent natural transformations of the corresponding telescope diagrams.  They therefore descend
to transformations
\[
    \id\longrightarrow R^{\mathrm{per}}L^{\mathrm{per}},
    \qquad
    L^{\mathrm{per}}R^{\mathrm{per}}\longrightarrow\id.
\]
The triangle identities hold stagewise and their coherent homotopies descend through the colimits,
so these transformations exhibit $L^{\mathrm{per}}\dashv R^{\mathrm{per}}$.

A compatible square of adjunctions gives a square of telescope diagrams and a transformation between
the two composites.  Periodization carries this data to the corresponding square and transformation.
Since formation of adjoint mates is characterized by the units and counits, the descended mate is the
mate of the descended transformation.  If the original transformation has a compatible inverse, the
inverse descends as well.
\end{proof}

\subsection{Compact generation under periodization}

\begin{lemma}[Compact closure of compact generators]
\label{lem:compact-closure-generators}
Let \(\mathcal C\) be a presentable \(\infty\)-category generated under small
colimits by a small family \(\mathcal G\) of compact objects.  Then
\(\mathcal C^\omega\) is the idempotent completion of the full subcategory
generated by \(\mathcal G\) under finite colimits.
\end{lemma}

\begin{proof}
Let
\[
    \mathcal D\subseteq\mathcal C
\]
be the full idempotent-complete subcategory generated by \(\mathcal G\) under
finite colimits and retracts.  Finite colimits of compact objects are compact,
and retracts of compact objects are compact, so
\[
    \mathcal D\subseteq\mathcal C^\omega.
\]

The inclusion
\[
    \mathcal D\hookrightarrow\mathcal C
\]
extends uniquely to a colimit-preserving functor
\[
    \Phi:
    \operatorname{Ind}(\mathcal D)
    \longrightarrow
    \mathcal C.
\]
Since both categories are presentable, \(\Phi\) admits a right adjoint
\[
    \Psi:
    \mathcal C
    \longrightarrow
    \operatorname{Ind}(\mathcal D).
\]

The functor \(\Phi\) carries every object of \(\mathcal D\) to a compact
object of \(\mathcal C\).  It follows that \(\Psi\) preserves filtered
colimits.  Indeed, for \(D\in\mathcal D\) and a filtered diagram
\(\{C_\alpha\}\) in \(\mathcal C\),
\[
\begin{aligned}
    \Map_{\operatorname{Ind}(\mathcal D)}
    \bigl(
      D,\Psi(\operatorname*{colim}_\alpha C_\alpha)
    \bigr)
    &\simeq
    \Map_{\mathcal C}
    \bigl(
      \Phi(D),\operatorname*{colim}_\alpha C_\alpha
    \bigr) \\
    &\simeq
    \operatorname*{colim}_\alpha
    \Map_{\mathcal C}
    \bigl(
      \Phi(D),C_\alpha
    \bigr) \\
    &\simeq
    \operatorname*{colim}_\alpha
    \Map_{\operatorname{Ind}(\mathcal D)}
    \bigl(
      D,\Psi(C_\alpha)
    \bigr).
\end{aligned}
\]
The objects of \(\mathcal D\) detect equivalences in
\(\operatorname{Ind}(\mathcal D)\), so \(\Psi\) preserves the filtered
colimit.

The unit
\[
    \eta:
    \id_{\operatorname{Ind}(\mathcal D)}
    \longrightarrow
    \Psi\Phi
\]
is an equivalence on \(\mathcal D\), because the inclusion
\(\mathcal D\hookrightarrow\mathcal C\) is fully faithful.  Both sides
preserve filtered colimits, and every object of
\(\operatorname{Ind}(\mathcal D)\) is a filtered colimit of objects of
\(\mathcal D\).  Hence \(\eta\) is an equivalence everywhere.  Thus
\(\Phi\) is fully faithful.

Let
\[
    \epsilon_C:
    \Phi\Psi(C)
    \longrightarrow
    C
\]
be the counit.  For every \(G\in\mathcal G\), adjunction and the unit
equivalence give
\[
\begin{aligned}
    \Map_{\mathcal C}
    \bigl(
      G,\Phi\Psi(C)
    \bigr)
    &\simeq
    \Map_{\operatorname{Ind}(\mathcal D)}
    \bigl(
      G,\Psi(C)
    \bigr) \\
    &\simeq
    \Map_{\mathcal C}(G,C).
\end{aligned}
\]
Under this identification, the map induced by \(\epsilon_C\) is the identity.
Since \(\mathcal G\) generates \(\mathcal C\), the counit \(\epsilon_C\) is an
equivalence.  Therefore
\[
    \Phi:
    \operatorname{Ind}(\mathcal D)
    \xrightarrow{\simeq}
    \mathcal C.
\]

The compact objects of \(\operatorname{Ind}(\mathcal D)\) are precisely the
objects of the idempotent-complete category \(\mathcal D\).  Hence
\[
    \mathcal C^\omega
    =
    \mathcal D,
\]
as asserted.
\end{proof}

\begin{proposition}[Compact generators after inversion]
\label{prop:compact-generation-periodization}
Let \(\mathcal C\) be compactly generated by a small family
\(\mathcal G\) of compact objects.  Suppose that \(X\) is symmetric and
compact and that tensoring with \(X\) preserves compact objects.  Then
\(\mathcal C[X^{-1}]\) is compactly generated by
\[
    X^{\otimes n}\otimes G,
    \qquad
    G\in\mathcal G,
    \quad
    n\in\mathbb Z,
\]
where the displayed objects mean
\[
    \operatorname{per}_X(X)^{\otimes n}
    \otimes
    \operatorname{per}_X(G)
\]
inside the periodized category.  Every object in this family is compact.
\end{proposition}

\begin{proof}
Let
\[
    \mathcal A:=\mathcal C^\omega.
\]
By Lemma~\ref{lem:compact-closure-generators}, \(\mathcal A\) is the
idempotent completion of the finite-colimit closure of \(\mathcal G\), and
there is an equivalence
\[
    \mathcal C
    \simeq
    \operatorname{Ind}(\mathcal A).
\]

The functor
\[
    X\otimes-:
    \mathcal C
    \longrightarrow
    \mathcal C
\]
preserves compact objects by assumption.  It therefore restricts to a
finite-colimit-preserving functor
\[
    X\otimes-:
    \mathcal A
    \longrightarrow
    \mathcal A.
\]
Form the small telescope
\[
    \mathcal A_X
    :=
    \operatorname{Idem}
    \operatorname*{colim}
    \left(
      \mathcal A
      \xrightarrow{X\otimes-}
      \mathcal A
      \xrightarrow{X\otimes-}
      \mathcal A
      \longrightarrow\cdots
    \right),
\]
where the colimit is taken among small \(\infty\)-categories with finite
colimits and finite-colimit-preserving functors, and
\(\operatorname{Idem}\) denotes idempotent completion.

The Ind-completion equivalence between small idempotent-complete
finite-colimit categories and compactly generated presentable categories with
compact-preserving left adjoints identifies
\[
    \operatorname{Ind}(\mathcal A_X)
\]
with the telescope
\[
    \operatorname*{colim}
    \left(
      \mathcal C
      \xrightarrow{X\otimes-}
      \mathcal C
      \xrightarrow{X\otimes-}
      \mathcal C
      \longrightarrow\cdots
    \right)
\]
in \(\Pr^L\).  By
Theorem~\ref{thm:periodization-telescope}, this telescope is canonically
equivalent to
\[
    \mathcal C[X^{-1}].
\]
Consequently,
\[
    \mathcal C[X^{-1}]
    \simeq
    \operatorname{Ind}(\mathcal A_X),
\]
and its compact objects are precisely the objects of \(\mathcal A_X\).

The images in \(\mathcal A_X\) of the objects of \(\mathcal A\) from all
telescope stages generate \(\mathcal A_X\) under finite colimits and
retracts.  Since \(\mathcal A\) is generated under finite colimits and
retracts by \(\mathcal G\), it is enough to use the images of the objects
\(G\in\mathcal G\) from all stages.

Transporting the image of \(G\) from the \(m\)-th stage to the zeroth stage
amounts to tensoring by
\[
    \operatorname{per}_X(X)^{-m}.
\]
Tensoring further by positive powers of the invertible object
\(\operatorname{per}_X(X)\) gives the family
\[
    \operatorname{per}_X(X)^{\otimes n}
    \otimes
    \operatorname{per}_X(G),
    \qquad
    n\in\mathbb Z.
\]
These objects are compact because they belong to \(\mathcal A_X\), and they
generate \(\mathcal C[X^{-1}]\) under small colimits.
\end{proof}


\subsection{Sequential \texorpdfstring{$\Omega_X$}{Omega-X}-spectra}

Let \(\mathcal C\) be a presentable pointed symmetric monoidal
\(\infty\)-category whose tensor product preserves small colimits separately
in each variable.  For \(X\in\mathcal C\), tensoring with \(X\) admits a
right adjoint
\[
    \Omega_X
    :=
    \underline{\operatorname{Hom}}(X,-):
    \mathcal C
    \longrightarrow
    \mathcal C.
\]

\begin{definition}[Sequential \(\Omega_X\)-spectra]
\label{def:X-prespectrum}
Define the \(\infty\)-category of sequential \(\Omega_X\)-spectra by
\[
    \Sp_X^\Omega(\mathcal C)
    :=
    \lim
    \left(
      \cdots
      \xrightarrow{\Omega_X}
      \mathcal C
      \xrightarrow{\Omega_X}
      \mathcal C
      \xrightarrow{\Omega_X}
      \mathcal C
    \right),
\]
where the limit is taken in \(\Pr^R\), equivalently in
\(\Cat_\infty\).

Thus an object of \(\Sp_X^\Omega(\mathcal C)\) consists of a sequence
\[
    E_0,E_1,E_2,\ldots
\]
together with equivalences
\[
    E_n
    \xrightarrow{\simeq}
    \underline{\operatorname{Hom}}
    \bigl(
      X,E_{n+1}
    \bigr)
\]
for every \(n\geq0\), coherently compatible with iteration.

For every \(n\geq0\), write
\[
    \operatorname{ev}_n:
    \Sp_X^\Omega(\mathcal C)
    \longrightarrow
    \mathcal C
\]
for evaluation at the \(n\)-th term.
\end{definition}

\begin{proposition}[Spectrum model for periodization]
\label{prop:spectrum-model-periodization}
Let \(\mathcal C\) be a presentable pointed symmetric monoidal
\(\infty\)-category whose tensor product preserves small colimits separately
in each variable, and let \(X\in\mathcal C\) be symmetric.

Then there is a canonical equivalence of presentable pointed
\(\infty\)-categories
\[
    \mathcal C[X^{-1}]
    \xrightarrow{\simeq}
    \Sp_X^\Omega(\mathcal C).
\]
Transporting the symmetric monoidal structure of
\(\mathcal C[X^{-1}]\) across this equivalence makes
\(\Sp_X^\Omega(\mathcal C)\) presentably symmetric monoidal.

Under this equivalence, the periodization functor
\[
    \operatorname{per}_X:
    \mathcal C
    \longrightarrow
    \mathcal C[X^{-1}]
\]
identifies with the left adjoint
\[
    \Sigma_X^\infty:
    \mathcal C
    \longrightarrow
    \Sp_X^\Omega(\mathcal C)
\]
of zeroth evaluation.  Thus there is an adjunction
\[
    \Sigma_X^\infty:
    \mathcal C
    \rightleftarrows
    \Sp_X^\Omega(\mathcal C):
    \operatorname{ev}_0.
\]

Let
\[
    F:
    \mathcal C
    \longrightarrow
    \mathcal D
\]
be a colimit-preserving strong symmetric monoidal functor, and put
\[
    U:=F(X).
\]
Then \(F\) induces a colimit-preserving strong symmetric monoidal functor
\[
    F_{\mathrm{sp}}:
    \Sp_X^\Omega(\mathcal C)
    \longrightarrow
    \Sp_U^\Omega(\mathcal D)
\]
with a canonical equivalence
\[
    F_{\mathrm{sp}}\Sigma_X^\infty
    \simeq
    \Sigma_U^\infty F.
\]
No preservation of internal Hom by \(F\) is required.

No stability assertion is made for
\(\mathcal C[X^{-1}]\) for a general symmetric object \(X\).
\end{proposition}

\begin{proof}
Consider the telescope
\[
    \operatorname{Tel}_X(\mathcal C)
    :=
    \operatorname*{colim}
    \left(
      \mathcal C
      \xrightarrow{X\otimes-}
      \mathcal C
      \xrightarrow{X\otimes-}
      \mathcal C
      \longrightarrow\cdots
    \right)
\]
in \(\Pr^L\).

Passing to right adjoints identifies the opposite of this colimit diagram
with the inverse diagram
\[
    \cdots
    \xrightarrow{\Omega_X}
    \mathcal C
    \xrightarrow{\Omega_X}
    \mathcal C
    \xrightarrow{\Omega_X}
    \mathcal C
\]
in \(\Pr^R\).  The equivalence
\[
    (\Pr^L)^{\op}
    \simeq
    \Pr^R
\]
therefore gives a canonical equivalence of presentable
\(\infty\)-categories
\[
    \operatorname{Tel}_X(\mathcal C)
    \xrightarrow{\simeq}
    \Sp_X^\Omega(\mathcal C).
\]
Under this equivalence, the canonical functor from the zeroth stage
\[
    \iota_0:
    \mathcal C
    \longrightarrow
    \operatorname{Tel}_X(\mathcal C)
\]
is left adjoint to
\[
    \operatorname{ev}_0:
    \Sp_X^\Omega(\mathcal C)
    \longrightarrow
    \mathcal C.
\]

Since \(X\) is symmetric,
Theorem~\ref{thm:periodization-telescope} gives a canonical equivalence
\[
    \operatorname{Tel}_X(\mathcal C)
    \xrightarrow{\simeq}
    \mathcal C[X^{-1}]
\]
under \(\mathcal C\).  Combining the two equivalences gives
\[
    \mathcal C[X^{-1}]
    \xrightarrow{\simeq}
    \Sp_X^\Omega(\mathcal C).
\]
Transport the presentably symmetric monoidal structure of
\(\mathcal C[X^{-1}]\) across this equivalence.

Under the combined equivalence, the canonical functor from the zeroth
telescope stage identifies with
\[
    \operatorname{per}_X.
\]
It therefore identifies with the left adjoint of
\(\operatorname{ev}_0\), which we denote by
\[
    \Sigma_X^\infty.
\]

Now let
\[
    F:
    \mathcal C
    \longrightarrow
    \mathcal D
\]
be colimit-preserving and strong symmetric monoidal, and put
\[
    U=F(X).
\]
Strong monoidality supplies a coherent equivalence
\[
    F(X\otimes C)
    \simeq
    U\otimes F(C).
\]
Thus \(F\) is \((X,U)\)-linear.
Proposition~\ref{prop:linear-functor-periodization} gives a
colimit-preserving functor
\[
    F^{\mathrm{per}}:
    \mathcal C[X^{-1}]
    \longrightarrow
    \mathcal D[U^{-1}]
\]
and a canonical equivalence
\[
    F^{\mathrm{per}}\operatorname{per}_X
    \simeq
    \operatorname{per}_U F.
\]
Because \(F\) is strong symmetric monoidal, the universal property of
periodization makes \(F^{\mathrm{per}}\) strong symmetric monoidal.

Transporting \(F^{\mathrm{per}}\) across the two spectrum-model
equivalences gives
\[
    F_{\mathrm{sp}}:
    \Sp_X^\Omega(\mathcal C)
    \longrightarrow
    \Sp_U^\Omega(\mathcal D)
\]
and
\[
    F_{\mathrm{sp}}\Sigma_X^\infty
    \simeq
    \Sigma_U^\infty F.
\]
This construction uses only the monoidal linearity of \(F\), not
preservation of either internal Hom functor.

Finally, formal inversion of an arbitrary symmetric object need not invert
the simplicial suspension object.  Hence no general stability conclusion
follows.
\end{proof}

\section{Stable brave new motivic spectra}
\label{sec:definition-SH}

\subsection{Definition and universal property}

\begin{definition}[Stable brave new motivic category]
\label{def:SH}
For a quasi-compact quasi-separated spectral scheme $S$, define
\[
    \SH(S)
    :=
    \HH_\bullet(S)[T_S^{-1}].
\]
The periodization functor is denoted
\[
    \Sigma_T^\infty:
    \HH_\bullet(S)
    \longrightarrow
    \SH(S).
\]
For an unpointed motivic space $F\in\HH(S)$, write
\[
    \Sigma_{T,+}^\infty F
    :=
    \Sigma_T^\infty(F_+).
\]
The motivic sphere spectrum is
\[
    \mathbf 1_S
    :=
    \Sigma_T^\infty\mathbb S^0_S.
\]
\end{definition}

\begin{proposition}[Universal property of $\SH(S)$]
\label{prop:SH-universal-property}
Let $\mathcal D$ be a presentably symmetric monoidal pointed $\infty$-category.  Restriction along
$\Sigma_T^\infty$ induces an equivalence
\[
\begin{aligned}
    \Fun^{L,\otimes}\bigl(\SH(S),\mathcal D\bigr)
    \simeq
    \bigl\{F\in\Fun^{L,\otimes}(\HH_\bullet(S),\mathcal D)
      \mid F(T_S)\text{ is invertible}\bigr\}.
\end{aligned}
\]
\end{proposition}

\begin{proof}
This is Theorem~\ref{thm:presentable-periodization} applied to
$\mathcal C=\HH_\bullet(S)$ and $X=T_S$.
\end{proof}

\begin{lemma}[Generators after periodization]
\label{lem:periodized-generators}
The category \(\SH(S)\) is generated under small colimits by the objects
\[
    T_S^{\otimes n}
    \otimes_S
    \Sigma_{T,S,+}^\infty\Mot_S(X),
    \qquad
    X\in\Sm^{\mathrm{fp}}_{/S},
    \quad
    n\in\mathbb Z.
\]
Equivalently, it is generated under small colimits by the images in the
periodization of the pointed smooth motives from every stage of the
\(T_S\)-telescope.
\end{lemma}

\begin{proof}
Chapter~4 shows that \(\HH_\bullet(S)\) is generated under small colimits by
the pointed smooth motives
\[
    \Mot_{S,+}(X),
    \qquad
    X\in\Sm^{\mathrm{fp}}_{/S}.
\]
By Proposition~\ref{prop:T-symmetric} and
Theorem~\ref{thm:periodization-telescope}, the underlying presentable category
of \(\SH(S)\) is the colimit
\[
    \HH_\bullet(S)
    \xrightarrow{T_S\wedge_S-}
    \HH_\bullet(S)
    \xrightarrow{T_S\wedge_S-}
    \HH_\bullet(S)
    \longrightarrow\cdots.
\]
The images of the generators from all stages generate this colimit.

Transporting the image of a pointed smooth motive from the \(m\)-th stage to
the zeroth stage amounts to tensoring its suspension spectrum by
\(T_S^{-m}\).  Closure under tensoring by positive powers of \(T_S\) gives
the displayed family indexed by all \(n\in\mathbb Z\).
\end{proof}

\begin{proposition}[Presentability and closed monoidality]
\label{prop:SH-presentable-closed}
The category $\SH(S)$ is presentable and symmetric monoidal.  Its tensor product
\[
    -\otimes_S-:\SH(S)\times\SH(S)\longrightarrow\SH(S)
\]
preserves small colimits separately in each variable and admits an internal Hom
\[
    \underline{\operatorname{Hom}}_S(-,-).
\]
The functor $\Sigma_T^\infty$ is colimit-preserving and strong symmetric monoidal.
\end{proposition}

\begin{proof}
Presentability, the symmetric monoidal structure, and strong monoidality of the periodization functor
are part of Theorem~\ref{thm:presentable-periodization}.  Separate preservation of colimits follows
because $\SH(S)$ is a commutative algebra object of $\Pr^L$.  A colimit-preserving tensor functor in
each variable on a presentable category admits a right adjoint in each variable, giving the internal
Hom.
\end{proof}

\subsection{Stability}

\begin{theorem}[Stability of brave new motivic spectra]
\label{thm:SH-stable}
The \(\infty\)-category \(\SH(S)\) is stable.
\end{theorem}

\begin{proof}
Let
\[
    S^1_S
    :=
    S^1\wedge\mathbb S^0_S
    \in
    \HH_\bullet(S)
\]
denote the simplicial suspension circle tensored with the pointed motivic
unit.  By construction,
\[
    \Sigma_T^\infty(T_S)
\]
is tensor-invertible in \(\SH(S)\).

Proposition~\ref{prop:T-circle-Gm} gives
\[
    T_S
    \simeq
    S^1_S\wedge_S\Gm{}_S
\]
in \(\HH_\bullet(S)\).  Since \(\Sigma_T^\infty\) is strong symmetric
monoidal,
\[
    \Sigma_T^\infty(T_S)
    \simeq
    \Sigma_T^\infty(S^1_S)
    \otimes_S
    \Sigma_T^\infty(\Gm{}_S).
\]
Lemma~\ref{lem:factors-invertible} therefore implies that
\[
    \Sigma_T^\infty(S^1_S)
\]
is tensor-invertible.

The ordinary suspension functor on the pointed presentable category
\(\SH(S)\) is tensoring with the simplicial circle:
\[
    \Sigma E
    \simeq
    \Sigma_T^\infty(S^1_S)\otimes_SE.
\]
It is therefore an equivalence, with inverse
\[
    E
    \longmapsto
    \Sigma_T^\infty(S^1_S)^{-1}\otimes_SE.
\]

A pointed presentable \(\infty\)-category is stable if and only if its
suspension functor is an equivalence.  The periodization
\[
    \SH(S)
    =
    \HH_\bullet(S)[T_S^{-1}]
\]
is pointed and presentable by
Theorem~\ref{thm:presentable-periodization}.  Hence \(\SH(S)\) is stable.
\end{proof}

\begin{corollary}[Exactness]
\label{cor:SH-exactness}
In $\SH(S)$, finite limits and finite colimits agree.  Every cofiber sequence is canonically a fiber
sequence, and conversely.
\end{corollary}

\begin{proof}
This is the defining exactness property of a stable $\infty$-category.
\end{proof}

\begin{notation}[Mapping spectra]
\label{not:mapping-spectra-SH}
The mapping spectrum in $\SH(S)$ is denoted
\[
    \operatorname{map}_S(E,F).
\]
Thus
\[
    \Omega^\infty\operatorname{map}_S(E,F)
    \simeq
    \Map_{\SH(S)}(E,F).
\]
\end{notation}

\subsection{The stable suspension and infinite-loop adjunction}

\begin{proposition}[Stable suspension adjunction]
\label{prop:SigmaOmegaT-adjunction}
The functor $\Sigma_T^\infty$ admits a right adjoint
\[
    \Omega_T^\infty:
    \SH(S)
    \longrightarrow
    \HH_\bullet(S).
\]
Thus there is an adjunction
\[
    \Sigma_T^\infty:
    \HH_\bullet(S)
    \rightleftarrows
    \SH(S):\Omega_T^\infty.
\]
The right adjoint preserves limits.
\end{proposition}

\begin{proof}
The functor $\Sigma_T^\infty$ preserves small colimits between presentable $\infty$-categories.
The adjoint functor theorem supplies a right adjoint.  Every right adjoint preserves limits.
\end{proof}

\begin{proposition}[Spectrum-object description]
\label{prop:SH-spectrum-object-description}
There is a canonical equivalence of presentable pointed
\(\infty\)-categories
\[
    \SH(S)
    \xrightarrow{\simeq}
    \Sp_{T_S}^\Omega
    \bigl(
      \HH_\bullet(S)
    \bigr).
\]
Equivalently,
\[
    \SH(S)
    \simeq
    \lim
    \left(
      \cdots
      \xrightarrow{\underline{\operatorname{Hom}}(T_S,-)}
      \HH_\bullet(S)
      \xrightarrow{\underline{\operatorname{Hom}}(T_S,-)}
      \HH_\bullet(S)
    \right).
\]

Under this presentation,
\[
    \Omega_T^\infty:
    \SH(S)
    \longrightarrow
    \HH_\bullet(S)
\]
is identified with zeroth evaluation.
\end{proposition}

\begin{proof}
The category \(\HH_\bullet(S)\) is presentable and pointed, and its smash
product preserves small colimits separately in each variable.
Proposition~\ref{prop:T-symmetric} shows that \(T_S\) is symmetric.
Proposition~\ref{prop:spectrum-model-periodization} therefore gives a
canonical equivalence
\[
\begin{aligned}
    \SH(S)
    &=
    \HH_\bullet(S)[T_S^{-1}] \\
    &\simeq
    \Sp_{T_S}^\Omega
    \bigl(
      \HH_\bullet(S)
    \bigr).
\end{aligned}
\]

Under this equivalence, the periodization functor
\[
    \Sigma_T^\infty:
    \HH_\bullet(S)
    \longrightarrow
    \SH(S)
\]
is the left adjoint of zeroth evaluation.  Its right adjoint is therefore
identified with
\[
    \operatorname{ev}_0:
    \Sp_{T_S}^\Omega
    \bigl(
      \HH_\bullet(S)
    \bigr)
    \longrightarrow
    \HH_\bullet(S).
\]
By uniqueness of right adjoints, this right adjoint is
\(\Omega_T^\infty\).
\end{proof}

\begin{remark}[Intrinsic suspension notation]
\label{rem:intrinsic-suspension-notation}
The fundamental stable suspension is the ordinary suspension $\Sigma$, equivalently tensoring with
$S^1$.  The additional motivic periodicity is encoded intrinsically by Thom twists.  Bigraded
notation will be introduced only after the virtual Thom construction, where its normalization is
unambiguous.
\end{remark}

\section{Stable base functoriality}
\label{sec:stable-base-functoriality}

\subsection{Inverse image}

Let $f:T\to S$ be a morphism of quasi-compact quasi-separated spectral schemes.  Pointed inverse
image is strong symmetric monoidal and Proposition~\ref{prop:thom-base-change-unstable} gives
\[
    f^*T_S
    \simeq
    T_T.
\]

\begin{proposition}[Stable inverse image]
\label{prop:stable-inverse-image}
There is a unique colimit-preserving strong symmetric monoidal functor
\[
    f^*:\SH(S)\longrightarrow\SH(T)
\]
with a canonical equivalence
\[
    f^*\Sigma_{T,S}^\infty
    \simeq
    \Sigma_{T,T}^\infty f^*.
\]
For composable morphisms
\[
    U\xrightarrow{g}T\xrightarrow{f}S,
\]
there are coherent equivalences
\[
    (fg)^*
    \simeq
    g^*f^*,
    \qquad
    (\id_S)^*
    \simeq
    \id_{\SH(S)}.
\]
\end{proposition}

\begin{proof}
The pointed inverse-image functor carries $T_S$ to $T_T$.  The universal property of
Proposition~\ref{prop:SH-universal-property} therefore gives the stable functor and its compatibility
with stabilization.  The unstable composition and identity equivalences preserve the stabilization
objects.  Their coherent descent through periodization gives the displayed equivalences.
\end{proof}

\begin{corollary}[The contravariant stable coefficient system]
\label{cor:SH-star-functor}
The assignments
\[
    S\longmapsto\SH(S),
    \qquad
    (f:T\to S)\longmapsto f^*
\]
define a functor
\[
    \SH^*:
    \bigl(\operatorname{Sch}^{\mathrm{sp}}_{\mathrm{qcqs}}\bigr)^{\mathrm{op}}
    \longrightarrow
    \CAlg(\Pr^L_{\mathrm{st}}).
\]
\end{corollary}

\begin{proof}
Proposition~\ref{prop:stable-inverse-image} supplies the objects, morphisms, and coherent
composition data.  Theorem~\ref{thm:SH-stable} places the values in stable presentable categories,
and stable inverse image is strong symmetric monoidal.
\end{proof}

\subsection{Direct image}

\begin{definition}[Stable direct image]
\label{def:stable-direct-image}
Define
\[
    f_*:\SH(T)\longrightarrow\SH(S)
\]
to be the right adjoint of stable inverse image.  Thus
\[
    f^*:\SH(S)
    \rightleftarrows
    \SH(T):f_*.
\]
\end{definition}

\begin{proposition}[Properties of stable direct image]
\label{prop:stable-direct-image-properties}
The functor $f_*$ preserves limits and is canonically lax symmetric monoidal.  For composable
morphisms $U\xrightarrow{g}T\xrightarrow{f}S$, there is a coherent equivalence
\[
    (fg)_*
    \simeq
    f_*g_*.
\]
There is a natural projection morphism
\[
    f_*E\otimes_SF
    \longrightarrow
    f_*\bigl(E\otimes_Tf^*F\bigr).
\]
\end{proposition}

\begin{proof}
Limit preservation is automatic for a right adjoint.  A right adjoint to a strong symmetric monoidal
functor is canonically lax symmetric monoidal; the displayed projection morphism is adjoint to the
composite
\[
\begin{aligned}
    f^*(f_*E\otimes_SF)
    &\simeq
    f^*f_*E\otimes_Tf^*F \\
    &\longrightarrow
    E\otimes_Tf^*F.
\end{aligned}
\]
The composition equivalence is the adjoint mate of the coherent equivalence
$(fg)^*\simeq g^*f^*$.
\end{proof}

\begin{proposition}[Compatibility with unstable direct image]
\label{prop:Omega-infty-direct-image}
There is a canonical equivalence
\[
    \Omega_{T,S}^\infty f_*
    \simeq
    f_*\Omega_{T,T}^\infty,
\]
where the functor on the right is pointed unstable direct image.
\end{proposition}

\begin{proof}
Proposition~\ref{prop:stable-inverse-image} gives a commutative square of left adjoints
\[
\begin{tikzcd}[column sep=large,row sep=large]
    \HH_\bullet(S) \arrow[r,"f^*"] \arrow[d,"\Sigma_{T,S}^\infty"']
        & \HH_\bullet(T) \arrow[d,"\Sigma_{T,T}^\infty"] \\
    \SH(S) \arrow[r,"f^*"'] & \SH(T).
\end{tikzcd}
\]
Taking right adjoints and using uniqueness of adjoint mates yields the asserted equivalence.
\end{proof}

\section{Stable smooth direct image}
\label{sec:stable-smooth-direct-image}

Let $p:X\to S$ be smooth and of finite presentation.

\subsection{$T$-linearity}

\begin{proposition}[$T$-linearity of pointed smooth direct image]
\label{prop:p-sharp-T-linear}
For every $F\in\HH_\bullet(X)$, there is a canonical equivalence
\[
    p_\sharp(T_X\wedge_XF)
    \simeq
    T_S\wedge_Sp_\sharp(F).
\]
These equivalences are coherent in $F$ and under composition of smooth morphisms.
\end{proposition}

\begin{proof}
By base change for Thom spaces,
\[
    p^*T_S
    \simeq
    T_X.
\]
The pointed smooth projection formula from Chapter~4 gives
\[
\begin{aligned}
    p_\sharp(T_X\wedge_XF)
    &\simeq
    p_\sharp(p^*T_S\wedge_XF) \\
    &\simeq
    T_S\wedge_Sp_\sharp(F).
\end{aligned}
\]
Coherence follows from the coherence of pullback, the pointed projection formula, and composition
of unstable smooth direct image.
\end{proof}

\subsection{Construction and adjunction}

\begin{theorem}[Stable smooth direct image]
\label{thm:stable-p-sharp}
There is a canonically induced colimit-preserving functor
\[
    p_\sharp:\SH(X)\longrightarrow\SH(S)
\]
with a canonical equivalence
\[
    p_\sharp\Sigma_{T,X}^\infty
    \simeq
    \Sigma_{T,S}^\infty p_\sharp.
\]
It is left adjoint to stable inverse image:
\[
    p_\sharp:\SH(X)
    \rightleftarrows
    \SH(S):p^*.
\]
\end{theorem}

\begin{proof}
Proposition~\ref{prop:p-sharp-T-linear} makes the unstable functor $p_\sharp$ $(T_X,T_S)$-linear.
Strong symmetric monoidality of $p^*$ and the equivalence $p^*T_S\simeq T_X$ make $p^*$
$(T_S,T_X)$-linear.  These two linearity equivalences are adjoint mates by the pointed smooth
projection formula.  Both functors preserve colimits.  Propositions~\ref{prop:T-symmetric},
\ref{prop:linear-functor-periodization}, and \ref{prop:periodized-adjunction} therefore supply the
stable adjunction and its compatibility with stabilization.
\end{proof}

\begin{proposition}[Smooth direct image on suspension spectra]
\label{prop:p-sharp-on-suspension-spectra}
Let $q:Y\to X$ be smooth and of finite presentation.  Then
\[
    p_\sharp\Sigma_{T,X,+}^\infty\Mot_X(Y)
    \simeq
    \Sigma_{T,S,+}^\infty\Mot_S(Y),
\]
where $Y$ on the right is regarded as a smooth spectral $S$-scheme by composition.
Equivalently,
\[
    p_\sharp\Sigma_{T,X,+}^\infty Y
    \simeq
    \Sigma_{T,S,+}^\infty Y.
\]
\end{proposition}

\begin{proof}
Before stabilization, smooth direct image carries $\Mot_{X,+}(Y)$ to $\Mot_{S,+}(Y)$.  Apply the
compatibility of Theorem~\ref{thm:stable-p-sharp} with $\Sigma_T^\infty$.
\end{proof}

\begin{proposition}[Coherent composition of stable smooth direct image]
\label{prop:stable-p-sharp-composition}
For smooth morphisms of finite presentation
\[
    X\xrightarrow{p}Y\xrightarrow{q}S,
\]
there is a canonical coherent equivalence
\[
    (qp)_\sharp
    \simeq
    q_\sharp p_\sharp.
\]
The identity smooth morphism induces the identity functor.
\end{proposition}

\begin{proof}
The corresponding equivalences hold for pointed unstable smooth direct image and are compatible
with the $T$-linearity data.  Proposition~\ref{prop:linear-functor-periodization} descends them to the
stable categories, retaining their coherence.
\end{proof}

\section{Stable smooth base change and projection}
\label{sec:stable-smooth-exchange}

\subsection{Smooth base change}

Consider a Cartesian square of quasi-compact quasi-separated spectral schemes
\[
\begin{tikzcd}[column sep=large,row sep=large]
    X_T \arrow[r,"g'"] \arrow[d,"p'"']
        & X \arrow[d,"p"] \\
    T \arrow[r,"g"'] & S,
\end{tikzcd}
\]
where $p$ is smooth and of finite presentation.  Then $p'$ is smooth and of finite presentation.

\begin{theorem}[Stable smooth base change]
\label{thm:stable-smooth-base-change}
The canonical Beck--Chevalley transformation
\[
    p'_\sharp(g')^*
    \longrightarrow
    g^*p_\sharp
\]
is an equivalence of functors
\[
    \SH(X)\longrightarrow\SH(T).
\]
\end{theorem}

\begin{proof}
The pointed unstable Beck--Chevalley transformation is an equivalence by Chapter~4.  Both composites
are $(T_X,T_T)$-linear: this follows from strong monoidality of inverse image and
Proposition~\ref{prop:p-sharp-T-linear}.  The transformation and its inverse are compatible with the linearity data and therefore descend
through periodization.  Their two composite homotopies descend as well, so the stable transformation
is an equivalence.
\end{proof}

\subsection{The smooth projection formula}

\begin{theorem}[Stable smooth projection formula]
\label{thm:stable-smooth-projection}
For \(E\in\SH(X)\) and \(F\in\SH(S)\), the canonical projection
transformation
\[
    p_\sharp\bigl(E\otimes_Xp^*F\bigr)
    \longrightarrow
    p_\sharp(E)\otimes_SF
\]
is an equivalence.
\end{theorem}

\begin{proof}
Let
\[
    \Phi_{E,F}:
    p_\sharp\bigl(E\otimes_Xp^*F\bigr)
    \longrightarrow
    p_\sharp(E)\otimes_SF
\]
denote the projection transformation.

By Lemma~\ref{lem:periodized-generators}, \(\SH(X)\) and \(\SH(S)\) are
generated under colimits by objects of the forms
\[
    E=
    T_X^{\otimes m}
    \otimes_X
    \Sigma_{T,X,+}^\infty\Mot_X(Y)
\]
and
\[
    F=
    T_S^{\otimes n}
    \otimes_S
    \Sigma_{T,S,+}^\infty\Mot_S(Z),
\]
where
\[
    Y\in\Sm^{\mathrm{fp}}_{/X},
    \qquad
    Z\in\Sm^{\mathrm{fp}}_{/S},
    \qquad
    m,n\in\mathbb Z.
\]

Using the \(T\)-linearity of \(p_\sharp\), the equivalence
\[
    p^*T_S\simeq T_X,
\]
and strong symmetric monoidality of \(p^*\), the transformation
\(\Phi_{E,F}\) reduces canonically to the case
\[
    m=n=0.
\]
In that case it is the stabilization of the pointed unstable projection
transformation
\[
\begin{aligned}
    p_\sharp\bigl(
      \Mot_{X,+}(Y)
      \wedge_X
      p^*\Mot_{S,+}(Z)
    \bigr)
    \longrightarrow
    p_\sharp\Mot_{X,+}(Y)
    \wedge_S
    \Mot_{S,+}(Z).
\end{aligned}
\]
Both sides are canonically identified with
\[
    \Mot_{S,+}(Y\times_S Z),
\]
and under these identifications the transformation is the identity.
Therefore \(\Phi_{E,F}\) is an equivalence on every pair of generators.

Fix a generator \(E\).  Both functors
\[
    F\longmapsto
    p_\sharp(E\otimes_Xp^*F)
\]
and
\[
    F\longmapsto
    p_\sharp(E)\otimes_SF
\]
preserve small colimits.  Hence \(\Phi_{E,F}\) is an equivalence for every
\(F\in\SH(S)\).

Now fix arbitrary \(F\).  Both functors
\[
    E\longmapsto
    p_\sharp(E\otimes_Xp^*F)
\]
and
\[
    E\longmapsto
    p_\sharp(E)\otimes_SF
\]
preserve small colimits.  Since the transformation is an equivalence on the
generators of \(\SH(X)\), it is an equivalence for every
\(E\in\SH(X)\).
\end{proof}

\begin{corollary}[Projection formula for pullback twists]
\label{cor:p-sharp-pullback-twist-preliminary}
Let $L\in\Pic(\SH(S))$.  Then
\[
    p_\sharp\bigl(E\otimes_Xp^*L\bigr)
    \simeq
    p_\sharp(E)\otimes_SL.
\]
\end{corollary}

\begin{proof}
Apply Theorem~\ref{thm:stable-smooth-projection} with $F=L$.
\end{proof}

\begin{remark}[Coherence]
\label{rem:smooth-exchange-coherence}
The smooth base-change and projection equivalences are induced from the coherent pointed unstable
transformations by periodization.  They are therefore compatible with iterated Cartesian squares,
composition of smooth morphisms, the monoidal constraints of inverse image, and adjoint mates.
These coherences will be used without further notation in the proper and exceptional theories.
\end{remark}

\section{Stable localization and recollement}
\label{sec:stable-localization}

Fix a complementary closed--open pair
\[
    Z\xrightarrow{i}S\xleftarrow{j}U,
\]
where $i$ is a closed immersion and $j$ is its quasi-compact open complement.

\subsection{Stabilization of the open and closed functors}

Since $j$ is smooth and of finite presentation, stable smooth direct image supplies
\[
    j_\sharp:\SH(U)
    \rightleftarrows
    \SH(S):j^*.
\]
Stable inverse image also supplies
\[
    j^*:\SH(S)
    \rightleftarrows
    \SH(U):j_*.
\]

The pointed unstable closed direct image
\[
    i_*:\HH_\bullet(Z)\longrightarrow\HH_\bullet(S)
\]
preserves colimits by the exactness theorem of Chapter~5.  The pointed closed projection formula
and $i^*T_S\simeq T_Z$ give
\[
\begin{aligned}
    i_*(T_Z\wedge_ZF)
    &\simeq
    i_*\bigl(i^*T_S\wedge_ZF\bigr) \\
    &\simeq
    T_S\wedge_Si_*F.
\end{aligned}
\]
Thus unstable closed direct image is $(T_Z,T_S)$-linear.

\begin{proposition}[Stable closed direct image]
\label{prop:stable-closed-direct-image}
There is a colimit-preserving functor
\[
    i_*:\SH(Z)\longrightarrow\SH(S)
\]
compatible with stabilization.  It is right adjoint to \(i^*\) and admits a
right adjoint \(i^!\):
\[
    i^*:\SH(S)
    \rightleftarrows
    \SH(Z):i_*
    \rightleftarrows
    \SH(S):i^!.
\]
Moreover, the functor \(i_*\) constructed by periodizing pointed unstable
closed direct image agrees canonically with the stable direct image of
Definition~\ref{def:stable-direct-image}.
\end{proposition}

\begin{proof}
The pointed unstable closed direct image
\[
    i_*:
    \HH_\bullet(Z)
    \longrightarrow
    \HH_\bullet(S)
\]
preserves all small colimits by the exactness theorem of Chapter~5.  The
pointed closed projection formula and the equivalence
\[
    i^*T_S\simeq T_Z
\]
give
\[
\begin{aligned}
    i_*(T_Z\wedge_ZF)
    &\simeq
    i_*\bigl(i^*T_S\wedge_ZF\bigr) \\
    &\simeq
    T_S\wedge_Si_*F.
\end{aligned}
\]
Thus pointed unstable closed direct image is
\((T_Z,T_S)\)-linear.  Proposition~\ref{prop:linear-functor-periodization}
therefore gives a colimit-preserving functor
\[
    i_*:
    \SH(Z)
    \longrightarrow
    \SH(S)
\]
compatible with stabilization.

Strong symmetric monoidality of pointed inverse image and the equivalence
\[
    i^*T_S\simeq T_Z
\]
make pointed \(i^*\) \((T_S,T_Z)\)-linear.  Its linearity equivalence and
that of pointed \(i_*\) are adjoint mates by the pointed closed projection
formula.  Both pointed functors preserve colimits.  Proposition
\ref{prop:periodized-adjunction} therefore gives an adjunction
\[
    i^*:
    \SH(S)
    \rightleftarrows
    \SH(Z):i_*.
\]

The stable functor \(i_*\) preserves small colimits between presentable
\(\infty\)-categories.  The adjoint functor theorem supplies a right adjoint
\[
    i^!:
    \SH(S)
    \longrightarrow
    \SH(Z).
\]

Definition~\ref{def:stable-direct-image} defines the stable direct image of
the morphism \(i\) as the right adjoint of stable inverse image
\[
    i^*:\SH(S)\longrightarrow\SH(Z).
\]
The periodized functor constructed above is also right adjoint to this same
functor.  Uniqueness of right adjoints therefore gives a canonical
equivalence between the two constructions.  This equivalence is compatible
with their units, counits, and stabilization comparison maps.
\end{proof}

\subsection{The localization triangles}

\begin{theorem}[Stable localization]
\label{thm:stable-localization}
For every \(E\in\SH(S)\), there is a natural exact triangle
\[
    j_\sharp j^*E
    \longrightarrow
    E
    \longrightarrow
    i_*i^*E.
\]
The first morphism is the counit of \(j_\sharp\dashv j^*\), and the second
is the unit of \(i^*\dashv i_*\).
\end{theorem}

\begin{proof}
For \(F\in\HH_\bullet(S)\), Chapter~5 gives a pointed cofiber sequence
\[
    j_\sharp j^*F
    \longrightarrow
    F
    \longrightarrow
    i_*i^*F.
\]
All three functors and both natural transformations are compatible with
stabilization.  Applying \(\Sigma_T^\infty\) therefore gives an exact
triangle
\[
    j_\sharp j^*\Sigma_T^\infty F
    \longrightarrow
    \Sigma_T^\infty F
    \longrightarrow
    i_*i^*\Sigma_T^\infty F.
\]

The localization transformation is \(T_S\)-linear.  Indeed,
\[
    j^*T_S\simeq T_U,
    \qquad
    i^*T_S\simeq T_Z,
\]
and the smooth and closed projection formulas give canonical equivalences
\[
\begin{aligned}
    j_\sharp j^*
    \bigl(
      T_S\otimes_SE
    \bigr)
    &\simeq
    T_S\otimes_Sj_\sharp j^*E, \\
    i_*i^*
    \bigl(
      T_S\otimes_SE
    \bigr)
    &\simeq
    T_S\otimes_Si_*i^*E.
\end{aligned}
\]
These equivalences identify the localization triangle for
\(T_S\otimes_SE\) with the tensor product by \(T_S\) of the localization
triangle for \(E\).

Since \(T_S\) is invertible in \(\SH(S)\), the same conclusion holds for
every integral power \(T_S^{\otimes n}\), including \(n<0\).  Consequently,
the asserted triangle is exact for every generator of the form
\[
    T_S^{\otimes n}
    \otimes_S
    \Sigma_{T,S}^\infty F,
    \qquad
    n\in\mathbb Z,
    \quad
    F\in\HH_\bullet(S).
\]

Let \(\mathcal L\subseteq\SH(S)\) be the full subcategory of objects \(E\)
for which the canonical morphism
\[
    \operatorname{cofib}
    \bigl(
      j_\sharp j^*E\longrightarrow E
    \bigr)
    \longrightarrow
    i_*i^*E
\]
is an equivalence.  The functors \(j_\sharp j^*\), \(\id\), and \(i_*i^*\)
preserve small colimits.  Since their source and target categories are
stable, they are exact.  It follows that \(\mathcal L\) is closed under small
colimits, suspensions, and desuspensions.

By Lemma~\ref{lem:periodized-generators}, the objects
\[
    T_S^{\otimes n}
    \otimes_S
    \Sigma_{T,S,+}^\infty\Mot_S(X),
    \qquad
    X\in\Sm^{\mathrm{fp}}_{/S},
    \quad
    n\in\mathbb Z,
\]
generate \(\SH(S)\) under small colimits.  They all belong to
\(\mathcal L\) by the preceding argument.  Therefore
\[
    \mathcal L=\SH(S),
\]
which proves the theorem.
\end{proof}

\begin{lemma}[Full faithfulness of stable open extension]
\label{lem:stable-j-sharp-fully-faithful}
The unit
\[
    \id_{\SH(U)}
    \longrightarrow
    j^*j_\sharp
\]
is an equivalence.  Consequently,
\[
    j_\sharp:\SH(U)\longrightarrow\SH(S)
\]
is fully faithful.
\end{lemma}

\begin{proof}
The corresponding pointed unstable unit
\[
    \id_{\HH_\bullet(U)}
    \longrightarrow
    j^*j_\sharp
\]
is an equivalence by the open-immersion formalism of Chapter~5.

The functors \(j_\sharp\) and \(j^*\) are compatible with the chosen
stabilization objects:
\[
    j^*T_S\simeq T_U,
\]
and the \(T\)-linearity of \(j_\sharp\) is the pointed smooth projection
formula.  The unstable unit is compatible with these linearity data.
Proposition~\ref{prop:periodized-adjunction} therefore carries it to an
equivalence
\[
    \id_{\SH(U)}
    \xrightarrow{\simeq}
    j^*j_\sharp.
\]
A left adjoint is fully faithful precisely when its unit is an equivalence.
\end{proof}

\begin{proposition}[Full faithfulness of stable open direct image]
\label{prop:stable-open-direct-image-fully-faithful}
The counit
\[
    j^*j_*\longrightarrow\id_{\SH(U)}
\]
is an equivalence.  Consequently, $j_*$ is fully faithful.
\end{proposition}

\begin{proof}
For $A,B\in\SH(U)$, the two open adjunctions and Lemma~\ref{lem:stable-j-sharp-fully-faithful} give natural
equivalences
\[
\begin{aligned}
    \Map_U(A,j^*j_*B)
    &\simeq
    \Map_S(j_\sharp A,j_*B) \\
    &\simeq
    \Map_U(j^*j_\sharp A,B) \\
    &\simeq
    \Map_U(A,B).
\end{aligned}
\]
Under this composite, the map induced by the counit $j^*j_*B\to B$ is the identity of
$\Map_U(A,B)$.  Yoneda therefore gives $j^*j_*B\simeq B$ naturally in $B$.
\end{proof}

\begin{theorem}[Stable co-localization]
\label{thm:stable-colocalization}
For every $E\in\SH(S)$, there is a natural exact triangle
\[
    i_*i^!E
    \longrightarrow
    E
    \longrightarrow
    j_*j^*E.
\]
The first morphism is the counit of $i_*\dashv i^!$ and the second is the unit of
$j^*\dashv j_*$.
\end{theorem}

\begin{proof}
The periodized counit
\[
    i^*i_*\longrightarrow\id_{\SH(Z)}
\]
is an equivalence, so $i_*$ is fully faithful.  The periodized open--closed orthogonality relation
gives
\[
    j^*i_*\simeq0.
\]
For $A\in\SH(Z)$ and $B\in\SH(U)$, adjunction and this orthogonality give
\[
\begin{aligned}
    \Map_Z(A,i^!j_*B)
    &\simeq
    \Map_S(i_*A,j_*B) \\
    &\simeq
    \Map_U(j^*i_*A,B) \\
    &\simeq *.
\end{aligned}
\]
Hence $i^!j_*\simeq0$.  Let $C(E)$ be the cofiber of $i_*i^!E\to E$.  Applying $i^!$ gives a
cofiber sequence
\[
    i^!i_*i^!E
    \longrightarrow
    i^!E
    \longrightarrow
    i^!C(E).
\]
The first morphism is an equivalence by full faithfulness, hence $i^!C(E)\simeq0$.

The unit map $E\to j_*j^*E$ annihilates $i_*i^!E$ by $j^*i_*\simeq0$, and therefore factors
canonically through a morphism
\[
    C(E)\longrightarrow j_*j^*E.
\]
After applying $j^*$, this morphism becomes the identity of $j^*E$.  After applying $i^!$, both
source and target vanish.  The pair $(i^!,j^*)$ is conservative: if $K$ is killed by both functors,
then stable localization gives $K\simeq i_*i^*K$, while
\[
    \Map_S(i_*i^*K,K)
    \simeq
    \Map_Z(i^*K,i^!K)
\]
is contractible, forcing the identity of $K$ to be null and hence $K\simeq0$.  Thus
$C(E)\to j_*j^*E$ is an equivalence, giving the asserted triangle.
\end{proof}

\subsection{The recollement package}

\begin{proposition}[Stable recollement]
\label{prop:stable-recollement}
The following assertions hold.
\begin{enumerate}
    \item The functors $i_*$, $j_\sharp$, and $j_*$ are fully faithful.

    \item There are canonical zero equivalences
    \[
        j^*i_*\simeq0,
        \qquad
        i^*j_\sharp\simeq0,
        \qquad
        i^!j_*\simeq0.
    \]

    \item An object $E\in\SH(S)$ lies in the essential image of $i_*$ if and only if
    $j^*E\simeq0$.

    \item The pair $(i^*,j^*)$ is conservative.
\end{enumerate}
Thus there is a recollement diagram
\[
    \SH(Z)
    \mathop{\longrightarrow}^{i_*}
    \mathop{\longleftarrow}_{i^*}
    \mathop{\longleftarrow}_{i^!}
    \SH(S)
    \mathop{\longrightarrow}^{j^*}
    \mathop{\longleftarrow}_{j_\sharp}
    \mathop{\longleftarrow}_{j_*}
    \SH(U).
\]
\end{proposition}

\begin{proof}
The unstable unit and counit equivalences proving full faithfulness of $i_*$ and $j_\sharp$ are
compatible with stabilization, so their stable counterparts are equivalences on the generating
$T$-shifts of suspension spectra and hence on all objects.  Full faithfulness of $j_*$ is
Proposition~\ref{prop:stable-open-direct-image-fully-faithful}.

The first two orthogonality equivalences descend directly from the pointed unstable theory.  The
third was proved in Theorem~\ref{thm:stable-colocalization}.  If $j^*E\simeq0$, stable localization gives
\[
    E\simeq i_*i^*E,
\]
so $E$ lies in the image of $i_*$.  The converse follows from $j^*i_*\simeq0$.  Finally, if both
$i^*E$ and $j^*E$ vanish, stable localization gives $E\simeq0$.
\end{proof}

\begin{corollary}[Support]
\label{cor:stable-support}
Let $\SH_Z(S)\subseteq\SH(S)$ be the full subcategory of objects $E$ with $j^*E\simeq0$.  Then
closed direct image induces an equivalence
\[
    i_*:\SH(Z)\xrightarrow{\simeq}\SH_Z(S).
\]
\end{corollary}

\begin{proof}
This is Proposition~\ref{prop:stable-recollement}(1) and (3).
\end{proof}

\subsection{Closed exchange formulas}

Consider a Cartesian square
\[
\begin{tikzcd}[column sep=large,row sep=large]
    Z_T \arrow[r,"g'"] \arrow[d,"i'"']
        & Z \arrow[d,"i"] \\
    T \arrow[r,"g"'] & S,
\end{tikzcd}
\]
where $i$ is closed.

\begin{theorem}[Stable closed base change]
\label{thm:stable-closed-base-change}
The canonical transformation
\[
    g^*i_*
    \longrightarrow
    i'_*(g')^*
\]
is an equivalence.
\end{theorem}

\begin{proof}
The pointed unstable closed base-change transformation is an equivalence by Chapter~5.  The four
functors in the square are compatible with the relevant Thom objects, and closed direct image is
$T$-linear.  Periodization therefore carries the unstable equivalence to the displayed stable
equivalence.
\end{proof}

\begin{theorem}[Stable closed projection formula]
\label{thm:stable-closed-projection}
For \(E\in\SH(Z)\) and \(F\in\SH(S)\), the canonical transformation
\[
    i_*E\otimes_SF
    \longrightarrow
    i_*\bigl(E\otimes_Zi^*F\bigr)
\]
is an equivalence.
\end{theorem}

\begin{proof}
Consider the pointed unstable bifunctors
\[
\begin{aligned}
    \Phi,\Psi:
    \HH_\bullet(Z)\times\HH_\bullet(S)
    &\longrightarrow
    \HH_\bullet(S), \\
    \Phi(E,F)
    &:=
    i_*E\wedge_SF, \\
    \Psi(E,F)
    &:=
    i_*\bigl(E\wedge_Zi^*F\bigr).
\end{aligned}
\]
Both bifunctors preserve small colimits separately in each variable.  For
\(\Phi\), this follows from colimit preservation of pointed closed direct
image and of the smash product.  For \(\Psi\), it follows additionally from
colimit preservation of \(i^*\).

The closed projection formula and the base-change equivalence
\[
    i^*T_S\simeq T_Z
\]
give coherent linearity equivalences in the first variable:
\[
\begin{aligned}
    \Phi(T_Z\wedge_ZE,F)
    &=
    i_*(T_Z\wedge_ZE)\wedge_SF \\
    &\simeq
    (T_S\wedge_Si_*E)\wedge_SF \\
    &\simeq
    T_S\wedge_S\Phi(E,F),
\end{aligned}
\]
and
\[
\begin{aligned}
    \Psi(T_Z\wedge_ZE,F)
    &=
    i_*\bigl(
      T_Z\wedge_ZE\wedge_Zi^*F
    \bigr) \\
    &\simeq
    T_S\wedge_S
    i_*\bigl(E\wedge_Zi^*F\bigr) \\
    &\simeq
    T_S\wedge_S\Psi(E,F).
\end{aligned}
\]

They are likewise linear in the second variable:
\[
\begin{aligned}
    \Phi(E,T_S\wedge_SF)
    &\simeq
    T_S\wedge_S\Phi(E,F),
\end{aligned}
\]
while
\[
\begin{aligned}
    \Psi(E,T_S\wedge_SF)
    &=
    i_*\bigl(
      E\wedge_Zi^*(T_S\wedge_SF)
    \bigr) \\
    &\simeq
    i_*\bigl(
      E\wedge_ZT_Z\wedge_Zi^*F
    \bigr) \\
    &\simeq
    T_S\wedge_S\Psi(E,F).
\end{aligned}
\]
The associativity and symmetry constraints of the smash product and the
coherence of the pointed closed projection formula give the required
commutation coherence between the two linearity structures.

Chapter~5 supplies a pointed unstable projection equivalence
\[
    \phi_{E,F}:
    \Phi(E,F)
    \xrightarrow{\simeq}
    \Psi(E,F).
\]
Its construction from the closed projection transformation makes it
compatible with the preceding multilinearity data.  Applying
Proposition~\ref{prop:multilinear-periodization} gives an equivalence of
periodized bifunctors
\[
    \Phi^{\mathrm{per}}
    \xrightarrow{\simeq}
    \Psi^{\mathrm{per}}.
\]

By construction of stable closed direct image, stable inverse image, and the
stable tensor products, these periodized bifunctors are
\[
    \Phi^{\mathrm{per}}(E,F)
    \simeq
    i_*E\otimes_SF
\]
and
\[
    \Psi^{\mathrm{per}}(E,F)
    \simeq
    i_*\bigl(E\otimes_Zi^*F\bigr).
\]
The periodized transformation is therefore the stable closed projection
transformation in the statement, and it is an equivalence.
\end{proof}

\begin{theorem}[Stable smooth--closed exchange]
\label{thm:stable-smooth-closed-exchange}
Let
\[
\begin{tikzcd}[column sep=large,row sep=large]
    Y \arrow[r,"k"] \arrow[d,"q"']
        & X \arrow[d,"p"] \\
    Z \arrow[r,"i"'] & S
\end{tikzcd}
\]
be Cartesian, where $p$ is smooth of finite presentation and $i$ is closed.  Then there is a
canonical equivalence
\[
    p_\sharp k_*
    \simeq
    i_*q_\sharp
\]
of functors $\SH(Y)\to\SH(S)$.
\end{theorem}

\begin{proof}
The corresponding pointed unstable exchange equivalence was proved in Chapter~5.  The functors
$p_\sharp$ and $q_\sharp$ are linear by the smooth projection formula, and $i_*$ and $k_*$ are
linear by the closed projection formula.  Hence the unstable equivalence descends through
periodization.
\end{proof}

\begin{proposition}[Internal-Hom exchange for a closed immersion]
\label{prop:closed-internal-hom}
For $F,G\in\SH(S)$, there is a canonical equivalence
\[
    i^!\underline{\operatorname{Hom}}_S(F,G)
    \simeq
    \underline{\operatorname{Hom}}_Z
    \bigl(i^*F,i^!G\bigr).
\]
Equivalently, for $E\in\SH(Z)$ and $G\in\SH(S)$,
\[
    \underline{\operatorname{Hom}}_S(i_*E,G)
    \simeq
    i_*\underline{\operatorname{Hom}}_Z(E,i^!G).
\]
\end{proposition}

\begin{proof}
For $K\in\SH(Z)$, the closed projection formula and the two adjunctions give
\[
\begin{aligned}
    \Map_Z\bigl(K,i^!\underline{\operatorname{Hom}}_S(F,G)\bigr)
    &\simeq
    \Map_S\bigl(i_*K,\underline{\operatorname{Hom}}_S(F,G)\bigr) \\
    &\simeq
    \Map_S(i_*K\otimes_SF,G) \\
    &\simeq
    \Map_S\bigl(i_*(K\otimes_Zi^*F),G\bigr) \\
    &\simeq
    \Map_Z(K\otimes_Zi^*F,i^!G) \\
    &\simeq
    \Map_Z\bigl(K,\underline{\operatorname{Hom}}_Z(i^*F,i^!G)\bigr).
\end{aligned}
\]
Yoneda gives the first equivalence.  For the second, map an arbitrary $F\in\SH(S)$ into both sides
and apply the first equivalence together with the closed projection formula.
\end{proof}

\section{Nisnevich descent for the stable coefficient system}
\label{sec:stable-category-descent}

The invertibility of a general Thom object will be proved locally.  We therefore record the
category-level descent statement needed for gluing inverses.

\subsection{Distinguished squares}

Let
\[
\begin{tikzcd}[column sep=large,row sep=large]
    W \arrow[r,"k"] \arrow[d,"q"']
        & V \arrow[d,"p"] \\
    U \arrow[r,"j"'] & X
\end{tikzcd}
\]
be a spectral Nisnevich distinguished square: $j$ is open, $p$ is \`etale, and $p$ induces an
equivalence over the closed complement of $U$.

\begin{definition}[The Nisnevich topology on spectral bases]
\label{def:base-Nisnevich-topology}
Let
\[
    \tau_{\mathrm{Nis}}^{\mathrm{base}}
\]
denote the Grothendieck topology on
\[
    \operatorname{Sch}^{\mathrm{sp}}_{\mathrm{qcqs}}
\]
generated by the empty covering family of the empty spectral scheme and by
the two-member covering families
\[
    \{U\longrightarrow X,\;V\longrightarrow X\}
\]
arising from Cartesian squares
\[
\begin{tikzcd}[column sep=large,row sep=large]
    W \arrow[r,"k"] \arrow[d,"q"']
        & V \arrow[d,"p"] \\
    U \arrow[r,"j"']
        & X
\end{tikzcd}
\]
such that:
\begin{enumerate}
    \item \(j\) is a quasi-compact open immersion;
    \item \(p\) is \'{e}tale and of finite presentation;
    \item if \(Z\hookrightarrow X\) is the complementary closed immersion,
    then the induced morphism
    \[
        V\times_X Z\longrightarrow Z
    \]
    is an equivalence.
\end{enumerate}
Such a square is called a Nisnevich distinguished square of spectral bases.
\end{definition}

\begin{lemma}[Pullbacks of presentable coefficient categories]
\label{lem:presentable-category-pullback}
Let
\[
    \mathcal A
    \xrightarrow{F}
    \mathcal C
    \xleftarrow{G}
    \mathcal B
\]
be colimit-preserving functors between presentable
\(\infty\)-categories.  Let
\[
    \mathcal P
    :=
    \mathcal A
    \times_{\mathcal C}
    \mathcal B
\]
be the \(\infty\)-category whose objects are triples
\[
    (A,B,\alpha),
    \qquad
    \alpha:F(A)\xrightarrow{\simeq}G(B).
\]

Then \(\mathcal P\) is presentable, its colimits are computed componentwise,
and it is the pullback of the displayed diagram in \(\Pr^L\).

If the three categories are presentably symmetric monoidal and \(F\) and
\(G\) are strong symmetric monoidal, then \(\mathcal P\) has the
componentwise presentably symmetric monoidal structure and is the pullback in
\(\CAlg(\Pr^L)\).  The same assertion holds in the stable and pointed
subcategories.
\end{lemma}

\begin{proof}
The category of triples is an accessible limit of accessible
\(\infty\)-categories and accessible functors, and is therefore accessible.
Because \(F\) and \(G\) preserve colimits, a colimit of a diagram of triples
is computed by taking the colimits of its two components and the induced
equivalence in \(\mathcal C\).  Thus \(\mathcal P\) admits all small
colimits, and accessibility implies that it is presentable.

For every presentable \(\infty\)-category \(\mathcal D\), a
colimit-preserving functor
\[
    \mathcal D\longrightarrow\mathcal P
\]
is exactly a pair of colimit-preserving functors
\[
    \mathcal D\longrightarrow\mathcal A,
    \qquad
    \mathcal D\longrightarrow\mathcal B
\]
together with an equivalence between their composites to
\(\mathcal C\).  Hence
\[
\begin{aligned}
    \Fun^L(\mathcal D,\mathcal P)
    \simeq
    \Fun^L(\mathcal D,\mathcal A)
    \times_{\Fun^L(\mathcal D,\mathcal C)}
    \Fun^L(\mathcal D,\mathcal B).
\end{aligned}
\]
This is the pullback universal property in \(\Pr^L\).

When \(F\) and \(G\) are strong symmetric monoidal, define
\[
    (A,B,\alpha)\otimes(A',B',\alpha')
    :=
    \bigl(
      A\otimes A',
      B\otimes B',
      \alpha\otimes\alpha'
    \bigr).
\]
The tensor product preserves colimits separately because the componentwise
tensor products do.  The same universal-property argument, now applied to
strong symmetric monoidal functors, gives the pullback in
\(\CAlg(\Pr^L)\).

Pointedness and stability are detected componentwise, so the final assertion
follows.
\end{proof}

\begin{proposition}[Stable Nisnevich excision of categories]
\label{prop:SH-Nis-square}
For every Nisnevich distinguished square of quasi-compact quasi-separated
spectral schemes
\[
\begin{tikzcd}[column sep=large,row sep=large]
    W \arrow[r,"k"] \arrow[d,"q"']
        & V \arrow[d,"p"] \\
    U \arrow[r,"j"']
        & X,
\end{tikzcd}
\]
restriction induces an equivalence of presentably symmetric monoidal stable
\(\infty\)-categories
\[
    \SH(X)
    \xrightarrow{\simeq}
    \SH(U)\times_{\SH(W)}\SH(V).
\]
\end{proposition}

\begin{proof}
Let
\[
    Z:=X\setminus U,
    \qquad
    Z_V:=V\setminus W,
\]
and write
\[
    i:Z\hookrightarrow X,
    \qquad
    i':Z_V\hookrightarrow V
\]
for the complementary closed immersions.  By the Nisnevich condition, the
induced morphism
\[
    e:Z_V\xrightarrow{\simeq}Z
\]
is an equivalence.

By Lemma~\ref{lem:presentable-category-pullback}, the pullback
\[
    \mathcal P
    :=
    \SH(U)\times_{\SH(W)}\SH(V)
\]
in \(\CAlg(\Pr^L_{\mathrm{st}})\) has underlying \(\infty\)-category the
category of triples
\[
    (A,B,\alpha),
    \qquad
    \alpha:q^*A\xrightarrow{\simeq}k^*B.
\]
It is stable and presentably symmetric monoidal, with colimits and tensor
products computed componentwise.

Let
\[
    J^*:
    \mathcal P
    \longrightarrow
    \SH(U)
\]
be the first projection.  It admits a left adjoint
\[
    J_\sharp:
    \SH(U)
    \longrightarrow
    \mathcal P
\]
defined by
\[
    J_\sharp(A)
    :=
    \bigl(
      A,
      k_\sharp q^*A,
      \eta_A
    \bigr),
\]
where
\[
    \eta_A:
    q^*A
    \xrightarrow{\simeq}
    k^*k_\sharp q^*A
\]
is the unit of
\[
    k_\sharp\dashv k^*.
\]
It is an equivalence because \(k_\sharp\) is fully faithful.

The functor \(J^*\) also admits a right adjoint
\[
    J_*:
    \SH(U)
    \longrightarrow
    \mathcal P
\]
defined by
\[
    J_*(A)
    :=
    \bigl(
      A,
      k_*q^*A,
      \epsilon_A^{-1}
    \bigr),
\]
where
\[
    \epsilon_A:
    k^*k_*q^*A
    \xrightarrow{\simeq}
    q^*A
\]
is the counit of
\[
    k^*\dashv k_*.
\]

Define
\[
    I_*:
    \SH(Z)
    \longrightarrow
    \mathcal P
\]
by
\[
    I_*(C)
    :=
    \bigl(
      0,
      i'_*e^*C,
      0
    \bigr).
\]
Since \(e\) is an equivalence, \(I_*\) has a left adjoint
\[
    I^*(A,B,\alpha)
    :=
    e_*i'^*B
\]
and a right adjoint
\[
    I^!(A,B,\alpha)
    :=
    e_*i'^!B.
\]

Let
\[
    P=(A,B,\alpha)\in\mathcal P.
\]
The stable localization triangle on \(V\),
\[
    k_\sharp k^*B
    \longrightarrow
    B
    \longrightarrow
    i'_*i'^*B,
\]
together with the gluing equivalence
\[
    \alpha:q^*A\xrightarrow{\simeq}k^*B,
\]
assembles with the triangle
\[
    A\xrightarrow{\id}A\longrightarrow0
\]
on \(U\) to a natural exact triangle in \(\mathcal P\):
\[
    J_\sharp J^*P
    \longrightarrow
    P
    \longrightarrow
    I_*I^*P.
\]
Thus \(\mathcal P\) is recollement-glued from \(\SH(U)\) and \(\SH(Z)\).

Restriction defines a colimit-preserving strong symmetric monoidal functor
\[
    R:
    \SH(X)
    \longrightarrow
    \mathcal P
\]
by
\[
    R(E)
    :=
    \bigl(
      j^*E,
      p^*E,
      \beta_E
    \bigr),
\]
where
\[
    \beta_E:
    q^*j^*E
    \xrightarrow{\simeq}
    k^*p^*E
\]
is the canonical Cartesian base-change equivalence.

By construction,
\[
    J^*R
    \simeq
    j^*.
\]
Stable smooth base change for the Cartesian square gives
\[
    Rj_\sharp
    \simeq
    J_\sharp.
\]
Indeed, the \(U\)-component is the unit equivalence
\[
    j^*j_\sharp\simeq\id,
\]
and the \(V\)-component is
\[
    p^*j_\sharp
    \simeq
    k_\sharp q^*.
\]

Stable closed base change gives
\[
    Ri_*
    \simeq
    I_*.
\]
It also gives
\[
\begin{aligned}
    I^*R
    &\simeq
    e_*i'^*p^* \\
    &\simeq
    e_*e^*i^* \\
    &\simeq
    i^*.
\end{aligned}
\]

Stable smooth--closed exchange for the Cartesian square
\[
\begin{tikzcd}[column sep=large,row sep=large]
    Z_V \arrow[r,"i'"] \arrow[d,"e"']
        & V \arrow[d,"p"] \\
    Z \arrow[r,"i"']
        & X
\end{tikzcd}
\]
gives
\[
    p_\sharp i'_*
    \simeq
    i_*e_\sharp.
\]
Taking right adjoints gives
\[
    i'^!p^*
    \simeq
    e^*i^!.
\]
Therefore
\[
\begin{aligned}
    I^!R
    &\simeq
    e_*i'^!p^* \\
    &\simeq
    e_*e^*i^! \\
    &\simeq
    i^!.
\end{aligned}
\]

We now prove that \(R\) is fully faithful.  Let
\[
    E,F\in\SH(X).
\]
Mapping the stable localization triangle
\[
    j_\sharp j^*E
    \longrightarrow
    E
    \longrightarrow
    i_*i^*E
\]
into \(F\) gives a fiber sequence of mapping spectra
\[
    \operatorname{map}_Z
    \bigl(
      i^*E,
      i^!F
    \bigr)
    \longrightarrow
    \operatorname{map}_X(E,F)
    \longrightarrow
    \operatorname{map}_U
    \bigl(
      j^*E,
      j^*F
    \bigr).
\]

The recollement triangle in \(\mathcal P\) gives a fiber sequence
\[
\begin{aligned}
    \operatorname{map}_Z
    \bigl(
      I^*R(E),
      I^!R(F)
    \bigr)
    \longrightarrow
    \operatorname{map}_{\mathcal P}
    \bigl(
      R(E),
      R(F)
    \bigr)
    \longrightarrow
    \operatorname{map}_U
    \bigl(
      J^*R(E),
      J^*R(F)
    \bigr).
\end{aligned}
\]
The preceding identifications identify the two outer terms.  Hence they
identify the middle terms:
\[
    \operatorname{map}_X(E,F)
    \xrightarrow{\simeq}
    \operatorname{map}_{\mathcal P}
    \bigl(
      R(E),
      R(F)
    \bigr).
\]
Thus \(R\) is fully faithful.

It remains to prove essential surjectivity.  Let \(P\in\mathcal P\).  Its
recollement triangle is
\[
    J_\sharp J^*P
    \longrightarrow
    P
    \longrightarrow
    I_*I^*P.
\]
Equivalently, \(P\) is the cofiber of a connecting morphism
\[
    I_*I^*P[-1]
    \longrightarrow
    J_\sharp J^*P.
\]
Both endpoint objects lie in the essential image of \(R\):
\[
    J_\sharp J^*P
    \simeq
    Rj_\sharp J^*P
\]
and
\[
    I_*I^*P
    \simeq
    Ri_*I^*P.
\]
Since \(R\) is fully faithful, the connecting morphism lifts through a
contractible space of choices to a morphism
\[
    i_*I^*P[-1]
    \longrightarrow
    j_\sharp J^*P
\]
in \(\SH(X)\).  Let \(E\) be its cofiber.  Exactness of \(R\) gives
\[
    R(E)\simeq P.
\]
Therefore \(R\) is essentially surjective.

Hence
\[
    R:
    \SH(X)
    \xrightarrow{\simeq}
    \SH(U)\times_{\SH(W)}\SH(V)
\]
is an equivalence.  It is strong symmetric monoidal because restriction is
strong symmetric monoidal and the tensor product in the pullback is computed
componentwise.
\end{proof}


\begin{lemma}[The Nisnevich cd-structure on spectral bases]
\label{lem:base-Nis-cd-structure}
The class of Nisnevich distinguished squares from
Definition~\ref{def:base-Nisnevich-topology} defines a complete and regular
cd-structure on
\[
    \operatorname{Sch}^{\mathrm{sp}}_{\mathrm{qcqs}}.
\]
The topology generated by this cd-structure and by the empty covering of the
empty spectral scheme is
\[
    \tau_{\mathrm{Nis}}^{\mathrm{base}}.
\]

Consequently, a space-valued presheaf
\[
    F:
    \bigl(
      \operatorname{Sch}^{\mathrm{sp}}_{\mathrm{qcqs}}
    \bigr)^{\op}
    \longrightarrow
    \mathcal S
\]
is a \(\tau_{\mathrm{Nis}}^{\mathrm{base}}\)-sheaf if and only if:
\begin{enumerate}
    \item \(F(\varnothing)\) is contractible;
    \item \(F\) carries every Nisnevich distinguished square to a Cartesian
    square of spaces.
\end{enumerate}
\end{lemma}

\begin{proof}
Nisnevich distinguished squares are stable under arbitrary base change:
open immersions, étale morphisms of finite presentation, complementary
closed immersions, and the condition of being an equivalence over the closed
complement are all stable under base change.

Classical truncation preserves Cartesian squares, open immersions,
complementary closed immersions, and étale morphisms.  By the classical
characterization of spectral étale morphisms established in Chapter~3, a
square of spectral bases is Nisnevich distinguished precisely when its
classical truncation is a classical Nisnevich distinguished square.

The completeness and regularity diagrams for the spectral cd-structure
therefore truncate to the corresponding diagrams for the classical
Nisnevich cd-structure.  Conversely, because the relevant open, étale, and
closed-complement constructions are recovered from their classical
truncations by spectral base change, the classical completeness and
regularity diagrams lift canonically to spectral schemes.

Thus the spectral distinguished squares form a complete and regular
cd-structure and generate the topology of
Definition~\ref{def:base-Nisnevich-topology}.  The final assertion is the
standard sheaf criterion for a complete regular cd-structure; see
\cite{VoevodskyCdStructures}.
\end{proof}

\begin{theorem}[Nisnevich descent of \(\SH\)]
\label{thm:SH-Nis-descent}
The contravariant functor
\[
    \SH^*:
    \bigl(
      \operatorname{Sch}^{\mathrm{sp}}_{\mathrm{qcqs}}
    \bigr)^{\op}
    \longrightarrow
    \CAlg(\Pr^L_{\mathrm{st}})
\]
satisfies descent for
\(\tau_{\mathrm{Nis}}^{\mathrm{base}}\).

In particular, for every finite Nisnevich covering family
\[
    U=\coprod_\alpha U_\alpha
    \longrightarrow
    X
\]
with \v{C}ech nerve \(U^\bullet\), restriction induces an equivalence
\[
    \SH(X)
    \xrightarrow{\simeq}
    \lim_{[n]\in\Delta}
    \SH(U^n)
\]
in \(\CAlg(\Pr^L_{\mathrm{st}})\).
\end{theorem}

\begin{proof}
The empty spectral scheme is a strict initial object of
\[
    \operatorname{Sch}^{\mathrm{sp}}_{\mathrm{qcqs}}.
\]
Its smooth site contains only the empty spectral scheme, and the empty
covering condition forces every motivic sheaf to take it to a contractible
space.  Therefore
\[
    \HH_\bullet(\varnothing)\simeq *,
    \qquad
    \SH(\varnothing)\simeq *.
\]

Fix
\[
    \mathcal D
    \in
    \CAlg(\Pr^L_{\mathrm{st}})
\]
and define a space-valued presheaf
\[
    \mathcal F_{\mathcal D}(X)
    :=
    \Map_{\CAlg(\Pr^L_{\mathrm{st}})}
    \bigl(
      \mathcal D,
      \SH(X)
    \bigr).
\]
The preceding calculation gives
\[
    \mathcal F_{\mathcal D}(\varnothing)
    \simeq
    *.
\]

By Proposition~\ref{prop:SH-Nis-square}, the functor \(\SH^*\) carries every
Nisnevich distinguished square to a pullback square in
\(\CAlg(\Pr^L_{\mathrm{st}})\).  Since the representable functor
\[
    \Map_{\CAlg(\Pr^L_{\mathrm{st}})}
    (\mathcal D,-)
\]
preserves limits, \(\mathcal F_{\mathcal D}\) carries every distinguished
square to a Cartesian square of spaces.

Lemma~\ref{lem:base-Nis-cd-structure} therefore implies that
\[
    \mathcal F_{\mathcal D}
\]
is a \(\tau_{\mathrm{Nis}}^{\mathrm{base}}\)-sheaf.

This conclusion holds for every test object
\[
    \mathcal D
    \in
    \CAlg(\Pr^L_{\mathrm{st}}).
\]
The Yoneda embedding
\[
    \CAlg(\Pr^L_{\mathrm{st}})
    \longrightarrow
    \Fun
    \bigl(
      \CAlg(\Pr^L_{\mathrm{st}})^{\op},
      \mathcal S
    \bigr)
\]
is fully faithful and preserves limits.  Hence the coefficient-system-valued
presheaf \(\SH^*\) itself satisfies
\(\tau_{\mathrm{Nis}}^{\mathrm{base}}\)-descent.

Applying the resulting descent equivalence to the covering family
\[
    U\longrightarrow X
\]
gives
\[
    \SH(X)
    \xrightarrow{\simeq}
    \lim_{[n]\in\Delta}
    \SH(U^n).
\]
\end{proof}

\subsection{Picard descent and local invertibility}

\begin{proposition}[Picard descent]
\label{prop:Pic-descent}
The contravariant functor
\[
    \Pic\SH:
    \bigl(
      \operatorname{Sch}^{\mathrm{sp}}_{\mathrm{qcqs}}
    \bigr)^{\op}
    \longrightarrow
    \CGrp(\mathcal S),
    \qquad
    X\longmapsto\Pic(\SH(X)),
\]
satisfies Nisnevich descent for
\(\tau_{\mathrm{Nis}}^{\mathrm{base}}\).
\end{proposition}

\begin{proof}
Let
\[
    \mathcal C
    \simeq
    \lim_\alpha\mathcal C_\alpha
\]
be a limit of symmetric monoidal \(\infty\)-categories.  An object of
\(\mathcal C\) is a compatible family
\[
    \{C_\alpha\}_\alpha.
\]
Such an object is tensor-invertible if and only if every \(C_\alpha\) is
tensor-invertible.  Indeed, compatible inverses in the components assemble
to an inverse in the limit, and the space of inverse data for an invertible
object is contractible.

Consequently, formation of the Picard \(\infty\)-groupoid preserves limits:
\[
    \Pic(\mathcal C)
    \simeq
    \lim_\alpha\Pic(\mathcal C_\alpha).
\]
Apply this observation to the Nisnevich descent equivalences of
Theorem~\ref{thm:SH-Nis-descent}.
\end{proof}

\begin{corollary}[Local criterion for invertibility]
\label{cor:invertibility-local}
Let $E\in\SH(X)$.  If there is a Nisnevich cover $\{U_\alpha\to X\}$ such that every restriction
$E|_{U_\alpha}$ is tensor-invertible, then $E$ is tensor-invertible.
\end{corollary}

\begin{proof}
The local inverses and their canonical compatibility equivalences define descent data in the Picard
$\infty$-groupoid.  Proposition~\ref{prop:Pic-descent} descends this data to an inverse of $E$ on
$X$.
\end{proof}

\section{Invertibility of Thom objects}
\label{sec:thom-invertibility}

\subsection{Stable Thom objects}

\begin{definition}[Stable Thom object]
\label{def:stable-thom-object}
For a finite locally free module $E$ on $X$, define
\[
    \operatorname{th}_X(E)
    :=
    \Sigma_T^\infty\Th_X(E)
    \in
    \SH(X).
\]
\end{definition}

\begin{remark}[The coordinate-module convention for Thom objects]
\label{rem:thom-dual-convention}
Under the standing convention
\[
    \mathbb V_X(E)
    =
    \operatorname{Spec}_X
    \operatorname{Sym}_{\mathcal O_X}(E),
\]
a morphism
\[
    T\longrightarrow\mathbb V_X(E)
\]
over \(X\) corresponds to an \(\mathcal O_T\)-linear morphism
\[
    E|_T\longrightarrow\mathcal O_T,
\]
and hence to a section of \(E^\vee|_T\).  Thus
\(\mathbb V_X(E)\) is the geometric vector bundle whose sheaf of sections is
\(E^\vee\).

Consequently,
\[
    \operatorname{th}_X(E)
    =
    \Sigma_T^\infty
    \left(
      \mathbb V_X(E)/
      \bigl(\mathbb V_X(E)\setminus X\bigr)
    \right)
\]
is the ordinary geometric Thom object of the vector bundle associated with
\(E^\vee\).

This convention agrees with the conormal convention of Chapter~5.  For a
quasi-smooth closed immersion
\[
    i:Z\hookrightarrow X,
\]
the geometric normal bundle is
\[
    \mathbf N_{Z/X}
    =
    \mathbb V_Z(N^*_{Z/X}),
\]
so its Thom space is indexed by the conormal module \(N^*_{Z/X}\).
Accordingly, all Thom twists in this chapter are indexed by the coordinate
module appearing in \(\mathbb V_X(-)\).
\end{remark}

\begin{theorem}[Stable homotopy purity]
\label{thm:stable-homotopy-purity}
Let
\[
    i:Z\lhook\joinrel\longrightarrow X
\]
be a smooth closed pair over \(S\), and let
\[
    z:Z\longrightarrow S
\]
be the structural morphism.  There is a canonical equivalence
\[
    \Sigma_{T,S}^\infty
    \bigl(
      X/(X\setminus Z)
    \bigr)
    \simeq
    z_\sharp
    \operatorname{th}_Z
    \bigl(
      N^*_{Z/X}
    \bigr)
\]
in \(\SH(S)\).

Equivalently, using the relative Thom-space notation,
\[
    \Sigma_{T,S}^\infty
    \Th_S
    \bigl(
      N^*_{Z/X}
    \bigr)
    \simeq
    z_\sharp
    \operatorname{th}_Z
    \bigl(
      N^*_{Z/X}
    \bigr).
\]

The equivalence is natural under arbitrary base change in \(S\) and under
\'{e}tale morphisms of smooth closed pairs.
\end{theorem}

\begin{proof}
Chapter~5 gives the unstable homotopy-purity equivalence
\[
    X/(X\setminus Z)
    \simeq
    \Th_S
    \bigl(
      N^*_{Z/X}
    \bigr)
\]
in \(\HH_\bullet(S)\).

By Proposition~\ref{prop:relative-thom-p-sharp}, there is a canonical
equivalence
\[
    z_\sharp
    \Th_Z
    \bigl(
      N^*_{Z/X}
    \bigr)
    \simeq
    \Th_S
    \bigl(
      N^*_{Z/X}
    \bigr).
\]
Apply
\[
    \Sigma_{T,S}^\infty
\]
to the unstable purity equivalence.  Compatibility of stable smooth direct
image with stabilization gives
\[
\begin{aligned}
    \Sigma_{T,S}^\infty
    z_\sharp
    \Th_Z
    \bigl(
      N^*_{Z/X}
    \bigr)
    &\simeq
    z_\sharp
    \Sigma_{T,Z}^\infty
    \Th_Z
    \bigl(
      N^*_{Z/X}
    \bigr) \\
    &=
    z_\sharp
    \operatorname{th}_Z
    \bigl(
      N^*_{Z/X}
    \bigr).
\end{aligned}
\]
Combining these equivalences proves the theorem.

Naturality follows from the base-change naturality of unstable homotopy
purity, base change for conormal modules and Thom spaces, and stable smooth
base change.
\end{proof}

Strong symmetric monoidality of stabilization and
Proposition~\ref{prop:unstable-thom-whitney} give
\[
    \operatorname{th}_X(E\oplus F)
    \simeq
    \operatorname{th}_X(E)\otimes_X\operatorname{th}_X(F).
\]
Moreover,
\[
    \operatorname{th}_X(\mathcal O_X^{\oplus n})
    \simeq
    T_X^{\otimes n},
\]
which is tensor-invertible by construction.

\begin{theorem}[Thom invertibility]
\label{thm:thom-invertibility}
Let \(X\) be a quasi-compact quasi-separated spectral scheme.  For every
finite locally free quasi-coherent \(\mathcal O_X\)-module \(E\), the stable
Thom object
\[
    \operatorname{th}_X(E)
\]
is tensor-invertible in \(\SH(X)\).
\end{theorem}

\begin{proof}
For every point \(x\in X\), finite local freeness of \(E\) gives an open
neighbourhood on which \(E\) is trivial.  Affine opens form a basis for the
topology of \(X\), so one may choose an affine open neighbourhood
\[
    U_x\longrightarrow X
\]
such that
\[
    E|_{U_x}
    \simeq
    \mathcal O_{U_x}^{\oplus r_x}
\]
for some integer \(r_x\geq0\).

The opens \(U_x\) cover \(X\).  Since \(X\) is quasi-compact, there is a
finite subcover
\[
    \{U_\alpha\longrightarrow X\}_{\alpha=1}^m.
\]
Each \(U_\alpha\) is affine and hence quasi-compact and quasi-separated, so
all coefficient categories and pullback functors used below lie within the
scope of Convention~\ref{conv:stable-coefficient-scope}.

For every \(\alpha\), choose an equivalence
\[
    E|_{U_\alpha}
    \simeq
    \mathcal O_{U_\alpha}^{\oplus r_\alpha}.
\]
Base change for Thom spaces and the trivial-bundle calculation give
\[
\begin{aligned}
    \operatorname{th}_X(E)|_{U_\alpha}
    &\simeq
    \operatorname{th}_{U_\alpha}
    \bigl(
      E|_{U_\alpha}
    \bigr) \\
    &\simeq
    \operatorname{th}_{U_\alpha}
    \bigl(
      \mathcal O_{U_\alpha}^{\oplus r_\alpha}
    \bigr) \\
    &\simeq
    T_{U_\alpha}^{\otimes r_\alpha}.
\end{aligned}
\]
The last object is tensor-invertible because
\(T_{U_\alpha}\) is tensor-invertible by the definition of
\(\SH(U_\alpha)\).

The finite Zariski covering family
\[
    \{U_\alpha\longrightarrow X\}_{\alpha=1}^m
\]
is, in particular, a Nisnevich covering family.  Hence
\(\operatorname{th}_X(E)\) is Nisnevich-locally tensor-invertible.
Corollary~\ref{cor:invertibility-local} therefore implies that
\[
    \operatorname{th}_X(E)
\]
is tensor-invertible in \(\SH(X)\).
\end{proof}

\begin{corollary}[The stable Thom functor]
\label{cor:stable-thom-functor}
The stable Thom construction refines to a symmetric monoidal functor
\[
    \operatorname{th}_X:
    \bigl(\Vect(X)^\simeq,\oplus\bigr)
    \longrightarrow
    \Pic(\SH(X)).
\]
It is natural under pullback:
\[
    f^*\operatorname{th}_X(E)
    \simeq
    \operatorname{th}_Y(f^*E).
\]
\end{corollary}

\begin{proof}
Corollary~\ref{cor:unstable-thom-functor} and strong symmetric monoidality of stabilization give a
symmetric monoidal functor to $\SH(X)$.  Theorem~\ref{thm:thom-invertibility} shows that its image
lies in the Picard subgroupoid.  Pullback compatibility follows from
Proposition~\ref{prop:thom-base-change-unstable}.
\end{proof}

\section{Vector-bundle and virtual Thom twists}
\label{sec:virtual-thom-twists}

\subsection{Vector-bundle suspension}

\begin{definition}[Vector-bundle suspension]
\label{def:vector-bundle-suspension}
Let $E$ be finite locally free on $X$.  Define the Thom suspension autoequivalence
\[
    \Sigma^E:\SH(X)\longrightarrow\SH(X)
\]
by
\[
    \Sigma^E(M)
    :=
    M\otimes_X\operatorname{th}_X(E).
\]
Define $\Sigma^{-E}$ by tensoring with the inverse of $\operatorname{th}_X(E)$.
\end{definition}

\begin{proposition}[Calculus of vector-bundle suspensions]
\label{prop:vector-bundle-suspension-calculus}
For finite locally free modules $E$ and $F$ on $X$, there are coherent equivalences
\[
    \Sigma^{E\oplus F}
    \simeq
    \Sigma^E\Sigma^F,
    \qquad
    \Sigma^0
    \simeq
    \id_{\SH(X)},
\]
\[
    \Sigma^E\Sigma^{-E}
    \simeq
    \id_{\SH(X)}
    \simeq
    \Sigma^{-E}\Sigma^E.
\]
For $f:Y\to X$, there is a coherent equivalence
\[
    f^*\Sigma^E
    \simeq
    \Sigma^{f^*E}f^*.
\]
\end{proposition}

\begin{proof}
The first two equivalences follow from symmetric monoidality of the stable Thom functor.  The next
two are the defining inverse equivalences in the Picard groupoid.  For pullback, use strong symmetric
monoidality of $f^*$ and the base-change equivalence
\[
    f^*\operatorname{th}_X(E)
    \simeq
    \operatorname{th}_Y(f^*E).
\]
All coherence is inherited from the symmetric monoidal Thom functor and coherent base
functoriality.
\end{proof}

\subsection{Additivity for extensions}

The extension from vector bundles to algebraic $K$-theory requires compatibility not only with direct
sums but also with exact sequences.  This compatibility follows canonically from stable
smooth--closed exchange.

\begin{proposition}[The Thom transformation]
\label{prop:thom-transformation}
Let $E$ be finite locally free on $X$, with projection
$\pi_E:\mathbb V_X(E)\to X$ and zero section
$s_E:X\to\mathbb V_X(E)$.  There is a canonical equivalence
\[
    \operatorname{th}_X(E)
    \simeq
    \pi_{E\sharp}s_{E*}\mathbf 1_X.
\]
More generally, for every $M\in\SH(X)$ there is a canonical equivalence
\[
    \pi_{E\sharp}s_{E*}M
    \simeq
    \operatorname{th}_X(E)\otimes_XM.
\]
\end{proposition}

\begin{proof}
Let $j_E:\mathbb V_X(E)\setminus X\to\mathbb V_X(E)$ be the open complement of the zero section.
Stable localization in $\SH(\mathbb V_X(E))$ gives an exact triangle
\[
    j_{E\sharp}\mathbf 1_{\mathbb V_X(E)\setminus X}
    \longrightarrow
    \mathbf 1_{\mathbb V_X(E)}
    \longrightarrow
    s_{E*}\mathbf 1_X.
\]
Applying $\pi_{E\sharp}$ identifies the cofiber of the first morphism with the stabilization of
\[
    \mathbb V_X(E)/\bigl(\mathbb V_X(E)\setminus X\bigr).
\]
This is $\operatorname{th}_X(E)$, proving the first equivalence.

For $M\in\SH(X)$, the closed projection formula and $s_E^*\pi_E^*\simeq\id$ give
\[
    s_{E*}M
    \simeq
    s_{E*}\mathbf 1_X\otimes_{\mathbb V_X(E)}\pi_E^*M.
\]
The smooth projection formula then gives
\[
\begin{aligned}
    \pi_{E\sharp}s_{E*}M
    &\simeq
    \pi_{E\sharp}\bigl(
      s_{E*}\mathbf 1_X\otimes\pi_E^*M
    \bigr) \\
    &\simeq
    \pi_{E\sharp}s_{E*}\mathbf 1_X\otimes_XM \\
    &\simeq
    \operatorname{th}_X(E)\otimes_XM.
\end{aligned}
\]
\end{proof}


\begin{definition}[Exact structure on finite locally free modules]
\label{def:exact-structure-Vect}
Let \(X\) be a quasi-compact quasi-separated spectral scheme.  The
\(\infty\)-category
\[
    \Vect(X)
\]
is the full subcategory of \(\QCoh(X)\) spanned by finite locally free
modules.

A sequence
\[
    E'
    \longrightarrow
    E
    \longrightarrow
    E''
\]
in \(\Vect(X)\) is exact if it is a cofiber sequence in the stable
\(\infty\)-category \(\QCoh(X)\).  Equivalently, Zariski-locally on \(X\),
it is equivalent to a split sequence
\[
    E'
    \longrightarrow
    E'\oplus E''
    \longrightarrow
    E''.
\]

These exact sequences make \(\Vect(X)\) an exact \(\infty\)-category.
Pullback along a morphism of spectral schemes is exact.
\end{definition}

\begin{proposition}[Thom additivity for exact sequences]
\label{prop:thom-exact-additivity}
For every exact sequence
\[
    0\longrightarrow E'
    \longrightarrow E
    \longrightarrow E''
    \longrightarrow0
\]
of finite locally free modules on \(X\), there is a canonical equivalence
\[
    \operatorname{th}_X(E)
    \simeq
    \operatorname{th}_X(E')
    \otimes_X
    \operatorname{th}_X(E'').
\]
This equivalence is natural under morphisms of exact sequences and is
compatible with direct sums of exact sequences.
\end{proposition}

\begin{proof}
The inclusion
\[
    E'\longrightarrow E
\]
induces a morphism of relative linear spaces
\[
    r:
    \mathbb V_X(E)
    \longrightarrow
    \mathbb V_X(E').
\]
We first verify that \(r\) is smooth and of finite presentation.  This is
Zariski-local on \(X\).  On a Zariski open over which the exact sequence
splits, there is an equivalence
\[
    E\simeq E'\oplus E'',
\]
under which \(r\) identifies with the projection
\[
    \mathbb V_X(E')
    \times_X
    \mathbb V_X(E'')
    \longrightarrow
    \mathbb V_X(E').
\]
This projection is the base change of
\[
    \mathbb V_X(E'')
    \longrightarrow
    X,
\]
which is smooth and of finite presentation.  Hence \(r\) is smooth and of
finite presentation globally.

Pulling \(r\) back along the zero section
\[
    s_{E'}:
    X
    \longrightarrow
    \mathbb V_X(E')
\]
gives a Cartesian square
\[
\begin{tikzcd}[column sep=large,row sep=large]
    \mathbb V_X(E'') \arrow[r,"k"] \arrow[d,"\pi_{E''}"']
        & \mathbb V_X(E) \arrow[d,"r"] \\
    X \arrow[r,"s_{E'}"']
        & \mathbb V_X(E').
\end{tikzcd}
\]
Indeed, relative spectrum converts this square into the pushout equivalence
\[
    \operatorname{Sym}(E)
    \otimes_{\operatorname{Sym}(E')}
    \mathcal O_X
    \simeq
    \operatorname{Sym}(E'').
\]
The pushout of
\[
    E'\longrightarrow E\longleftarrow0
\]
in quasi-coherent modules is \(E''\), and the symmetric-algebra functor
preserves colimits.

The zero section of \(E\) factors as
\[
    s_E=ks_{E''},
\]
and the structure morphism factors as
\[
    \pi_E=\pi_{E'}r.
\]
Proposition~\ref{prop:thom-transformation} therefore gives
\[
\begin{aligned}
    \operatorname{th}_X(E)
    &\simeq
    \pi_{E\sharp}s_{E*}\mathbf1_X \\
    &\simeq
    \pi_{E'\sharp}
    r_\sharp
    k_*
    s_{E''*}\mathbf1_X.
\end{aligned}
\]
Stable smooth--closed exchange for the displayed Cartesian square gives
\[
    r_\sharp k_*
    \simeq
    s_{E'*}\pi_{E''\sharp}.
\]
Consequently,
\[
\begin{aligned}
    \operatorname{th}_X(E)
    &\simeq
    \pi_{E'\sharp}
    s_{E'*}
    \pi_{E''\sharp}
    s_{E''*}\mathbf1_X \\
    &\simeq
    \pi_{E'\sharp}
    s_{E'*}
    \operatorname{th}_X(E'') \\
    &\simeq
    \operatorname{th}_X(E')
    \otimes_X
    \operatorname{th}_X(E'').
\end{aligned}
\]
The final equivalence is Proposition~\ref{prop:thom-transformation} applied
to the object \(\operatorname{th}_X(E'')\).

A morphism of exact sequences induces a morphism between the corresponding
Cartesian squares of relative linear spaces.  Naturality of relative
spectrum, zero sections, smooth--closed exchange, and the smooth and closed
projection transformations gives naturality of the additivity equivalence.

For two exact sequences, the Cartesian square associated with their direct
sum is the product of the two corresponding Cartesian squares.  Compatibility
of smooth--closed exchange with products and the symmetric monoidal
projection formulas identifies the additivity equivalence for the direct sum
with the tensor product of the two additivity equivalences.
\end{proof}

\begin{proposition}[Coherent Thom additivity]
\label{prop:coherent-thom-additivity}
Let
\[
    S_\bullet\Vect(X)
\]
be the Waldhausen \(S\)-construction of the exact \(\infty\)-category
\(\Vect(X)\) from
Definition~\ref{def:exact-structure-Vect}.

For every \(n\geq0\), define
\[
    \delta_n:
    \bigl(
      S_n\Vect(X)
    \bigr)^\simeq
    \longrightarrow
    \Pic(\SH(X))^n
\]
by
\[
\begin{aligned}
    \delta_n
    \bigl(
      0=E_0\to E_1\to\cdots\to E_n
    \bigr)
    :=
    \bigl(
      \operatorname{th}_X(E_1/E_0),
      \ldots,
      \operatorname{th}_X(E_n/E_{n-1})
    \bigr).
\end{aligned}
\]

The maps \(\delta_n\) assemble into a simplicial map of
\(E_\infty\)-spaces
\[
    \delta_\bullet:
    \bigl(
      S_\bullet\Vect(X)
    \bigr)^\simeq
    \longrightarrow
    B_\bullet\Pic(\SH(X)).
\]
\end{proposition}

\begin{proof}
Let
\[
    0=E_0\longrightarrow E_1\longrightarrow\cdots\longrightarrow E_n
\]
be an object of \(S_n\Vect(X)\).  For \(0\leq a<b\leq n\), put
\[
    Q_{a,b}:=E_b/E_a.
\]
For every triple
\[
    a<b<c,
\]
there is a canonical exact sequence
\[
    0
    \longrightarrow
    Q_{a,b}
    \longrightarrow
    Q_{a,c}
    \longrightarrow
    Q_{b,c}
    \longrightarrow
    0.
\]
Proposition~\ref{prop:thom-exact-additivity} gives a natural equivalence
\[
    \operatorname{th}_X(Q_{a,c})
    \xrightarrow{\simeq}
    \operatorname{th}_X(Q_{a,b})
    \otimes_X
    \operatorname{th}_X(Q_{b,c}).
\]
Let
\[
\begin{aligned}
    \mu_{a,b,c}:
    \operatorname{th}_X(Q_{a,b})
    \otimes_X
    \operatorname{th}_X(Q_{b,c})
    \xrightarrow{\simeq}
    \operatorname{th}_X(Q_{a,c})
\end{aligned}
\]
denote its inverse.  This direction agrees with multiplication in the bar
construction.

For four indices
\[
    a<b<c<d,
\]
there are two composites
\[
\begin{aligned}
    &
    \operatorname{th}_X(Q_{a,b})
    \otimes
    \operatorname{th}_X(Q_{b,c})
    \otimes
    \operatorname{th}_X(Q_{c,d}) \\
    &\hspace{4cm}
    \longrightarrow
    \operatorname{th}_X(Q_{a,d}),
\end{aligned}
\]
obtained by first multiplying the first two factors or the last two factors.

The additivity equivalence of
Proposition~\ref{prop:thom-exact-additivity} is constructed from the
Cartesian square of relative linear spaces associated with a short exact
sequence, together with smooth--closed exchange and the smooth and closed
projection formulas.  For the filtration
\[
    E_a\subseteq E_b\subseteq E_c\subseteq E_d,
\]
the two composites above arise from the two possible pastings of the same
iterated Cartesian diagram.  The Beck--Chevalley pasting law identifies the
two exchange transformations, and compatibility of the projection formulas
with composition identifies the resulting tensor-product maps.  Thus the
two composites agree through a canonical homotopy.

More generally, for every subdivision
\[
    a=a_0<a_1<\cdots<a_r=b,
\]
the iterated multiplication
\[
\begin{aligned}
    \bigotimes_{\nu=1}^r
    \operatorname{th}_X
    \bigl(
      Q_{a_{\nu-1},a_\nu}
    \bigr)
    \longrightarrow
    \operatorname{th}_X(Q_{a,b})
\end{aligned}
\]
is obtained by pasting the same iterated diagram of relative linear spaces.
Higher Beck--Chevalley pasting coherence identifies all parenthesizations
and all refinements through a contractible space of choices.  These are
exactly the associativity coherences required for a map into the bar
construction.

We now check the simplicial operators.

An inner face map removes an intermediate term \(E_a\).  It replaces the
two consecutive quotients
\[
    Q_{a-1,a},
    \qquad
    Q_{a,a+1}
\]
by
\[
    Q_{a-1,a+1}.
\]
The corresponding face map in
\[
    B_\bullet\Pic(\SH(X))
\]
multiplies the two consecutive entries.  Compatibility is supplied by
\[
    \mu_{a-1,a,a+1}.
\]

The first and last face maps discard the first or last quotient.  This
agrees directly with the definition of \(\delta_n\).

A degeneracy repeats one term of the filtration and therefore inserts the
zero quotient.  Since
\[
    \operatorname{th}_X(0)
    \simeq
    \mathbf1_X,
\]
this agrees with insertion of the unit in the bar construction.

The coherent compatibility for arbitrary subdivisions proves all simplicial
identities, including their higher homotopies.  Therefore the maps
\(\delta_n\) assemble into a simplicial map
\[
    \delta_\bullet:
    \bigl(
      S_\bullet\Vect(X)
    \bigr)^\simeq
    \longrightarrow
    B_\bullet\Pic(\SH(X)).
\]

Finally, direct sum gives
\[
    \bigl(
      S_\bullet\Vect(X)
    \bigr)^\simeq
\]
a simplicial \(E_\infty\)-space.  The stable Thom functor is symmetric
monoidal:
\[
    \operatorname{th}_X(E\oplus F)
    \simeq
    \operatorname{th}_X(E)
    \otimes_X
    \operatorname{th}_X(F).
\]
The direct-sum compatibility in
Proposition~\ref{prop:thom-exact-additivity} shows that every
\(\delta_n\) is a map of \(E_\infty\)-spaces.  Hence
\(\delta_\bullet\) is a simplicial map of \(E_\infty\)-spaces.
\end{proof}

\subsection{The motivic $J$-homomorphism}

\begin{definition}[Vector-bundle algebraic \(K\)-theory]
\label{def:vector-bundle-K-theory}
Let
\[
    \Vect(X)
\]
denote the exact \(\infty\)-category of finite locally free
quasi-coherent \(\mathcal O_X\)-modules.  Define
\[
    K(X)
    :=
    K\bigl(\Vect(X)\bigr)
\]
to be its algebraic \(K\)-theory space, constructed by the Waldhausen
\(S\)-construction.  Direct sum makes \(K(X)\) a grouplike
\(E_\infty\)-space.

Thus \(K(X)\) in this chapter means vector-bundle algebraic \(K\)-theory.
No identification with the algebraic \(K\)-theory of all perfect complexes is
used or asserted.
\end{definition}

\begin{theorem}[The motivic \(J\)-homomorphism]
\label{thm:motivic-J}
The stable Thom functor and coherent Thom additivity induce a natural map of
grouplike \(E_\infty\)-spaces
\[
    J_X:
    K(X)
    \longrightarrow
    \Pic(\SH(X)).
\]
For every finite locally free module \(E\),
\[
    J_X([E])
    \simeq
    \operatorname{th}_X(E).
\]

For every morphism \(f:Y\to X\), the square
\[
\begin{tikzcd}[column sep=large,row sep=large]
    K(X) \arrow[r,"J_X"] \arrow[d,"f^*"']
        & \Pic(\SH(X)) \arrow[d,"f^*"] \\
    K(Y) \arrow[r,"J_Y"']
        & \Pic(\SH(Y))
\end{tikzcd}
\]
commutes through a canonical coherent homotopy.
\end{theorem}

\begin{proof}
Proposition~\ref{prop:coherent-thom-additivity} gives a simplicial map of
\(E_\infty\)-spaces
\[
    \delta_\bullet:
    \bigl(
      S_\bullet\Vect(X)
    \bigr)^\simeq
    \longrightarrow
    B_\bullet\Pic(\SH(X)).
\]
Geometric realization gives a map of \(E_\infty\)-spaces
\[
    \left|
      \bigl(
        S_\bullet\Vect(X)
      \bigr)^\simeq
    \right|
    \longrightarrow
    \left|
      B_\bullet\Pic(\SH(X))
    \right|.
\]
Looping gives
\[
\begin{aligned}
    J_X:
    K(X)
    &=
    \Omega
    \left|
      \bigl(
        S_\bullet\Vect(X)
      \bigr)^\simeq
    \right| \\
    &\longrightarrow
    \Omega
    \left|
      B_\bullet\Pic(\SH(X))
    \right| \\
    &\simeq
    \Pic(\SH(X)).
\end{aligned}
\]
Because \(\delta_\bullet\) is a simplicial map of \(E_\infty\)-spaces,
\(J_X\) is a map of grouplike \(E_\infty\)-spaces.

The one-step filtration
\[
    0\longrightarrow E
\]
represents the class of \(E\) and is carried by \(\delta_1\) to
\[
    \operatorname{th}_X(E).
\]
Hence
\[
    J_X([E])
    \simeq
    \operatorname{th}_X(E).
\]

Let \(f:Y\to X\).  Pullback preserves finite locally free modules,
cofiber sequences of finite locally free modules, quotients in
filtrations, and direct sums.  Base change for Thom objects gives
\[
    f^*\operatorname{th}_X(E)
    \simeq
    \operatorname{th}_Y(f^*E).
\]
Naturality of smooth--closed exchange, zero sections, relative linear
spaces, and the projection formulas shows that these equivalences commute
with the multiplication maps
\[
    \mu_{a,b,c}
\]
used in the construction of \(\delta_\bullet\).

Thus pullback commutes coherently with the complete simplicial
\(E_\infty\)-map
\[
    \delta_\bullet.
\]
Geometric realization and looping produce the asserted coherent homotopy
between
\[
    f^*J_X
    \qquad\text{and}\qquad
    J_Yf^*.
\]
\end{proof}

\begin{definition}[Virtual Thom suspension]
\label{def:virtual-thom-suspension}
For $\xi\in K(X)$, write
\[
    \operatorname{th}_X(\xi)
    :=
    J_X(\xi)
    \in
    \Pic(\SH(X)).
\]
Define
\[
    \Sigma^\xi(M)
    :=
    M\otimes_X\operatorname{th}_X(\xi).
\]
Given a specified path
\[
    \gamma:\xi\simeq[E]-[F]
\]
in $K(X)$, applying $J_X$ to $\gamma$ induces an equivalence
\[
    \Sigma^\xi
    \simeq
    \Sigma^E\Sigma^{-F}.
\]
\end{definition}

\begin{proposition}[Virtual-twist calculus]
\label{prop:virtual-twist-calculus}
For $\xi,\eta\in K(X)$, there are coherent equivalences
\[
    \Sigma^{\xi+\eta}
    \simeq
    \Sigma^\xi\Sigma^\eta,
    \qquad
    \Sigma^0\simeq\id,
    \qquad
    \Sigma^{-\xi}\simeq(\Sigma^\xi)^{-1}.
\]
For $f:Y\to X$,
\[
    f^*\Sigma^\xi
    \simeq
    \Sigma^{f^*\xi}f^*.
\]
The construction depends only on the point $\xi$ of the $K$-theory space and is compatible with
homotopies in $K(X)$.
\end{proposition}

\begin{proof}
These statements are the functorial consequences of the map of grouplike $E_\infty$-spaces
$J_X$ and its pullback naturality.  A path in $K(X)$ is carried to the corresponding equivalence in
the Picard $\infty$-groupoid, so the construction is compatible with all specified identifications of
virtual classes.
\end{proof}

\begin{proposition}[Smooth projection with virtual twists]
\label{prop:p-sharp-virtual-twist}
Let $p:X\to S$ be smooth of finite presentation, let $E\in\SH(X)$, and let $\xi\in K(S)$.  Then
there is a canonical equivalence
\[
    p_\sharp\bigl(\Sigma^{p^*\xi}E\bigr)
    \simeq
    \Sigma^\xi p_\sharp(E).
\]
\end{proposition}

\begin{proof}
By pullback compatibility,
\[
    \Sigma^{p^*\xi}E
    \simeq
    E\otimes_Xp^*\operatorname{th}_S(\xi).
\]
Apply the stable smooth projection formula.
\end{proof}

\begin{notation}[Bigraded suspension]
\label{not:bigraded-suspension}
For integers \(p,q\in\mathbb Z\), define
\[
    \Sigma^{p,q}
    :=
    \Sigma^{p-2q}
    \Sigma^{q\mathcal O_S}.
\]
Thus
\[
    \Sigma^{2,1}
    \simeq
    \Sigma^{\mathcal O_S},
\]
and tensoring with \(T_S\) is the suspension
\[
    \Sigma^{2,1}.
\]
Virtual-bundle notation remains the intrinsic notation and will be used
whenever a nontrivial bundle is involved.
\end{notation}

\section{Compact generation}
\label{sec:SH-compact-generation}

\subsection{Compactness of the stabilization object}



\begin{proposition}[Compactness of \(T_S\)]
\label{prop:T-compact}
The object
\[
    T_S\in\HH_\bullet(S)
\]
is compact.  Tensoring with \(T_S\) preserves compact objects.
\end{proposition}

\begin{proof}
Let
\[
    U:\HH_\bullet(S)\longrightarrow\HH(S)
\]
be the forgetful functor.  Since \(\HH_\bullet(S)\) is the
\(\infty\)-category of pointed objects in \(\HH(S)\), colimits are computed
on underlying motivic spaces.  In particular, \(U\) preserves filtered
colimits.

For every smooth spectral \(S\)-scheme \(Y\) of finite presentation and every
\(F\in\HH_\bullet(S)\), the free-pointing adjunction gives
\[
    \Map_{\HH_\bullet(S)}
    \bigl(\Mot_{S,+}(Y),F\bigr)
    \simeq
    \Map_{\HH(S)}
    \bigl(\Mot_S(Y),U(F)\bigr).
\]
The object \(\Mot_S(Y)\) is compact in \(\HH(S)\) by Chapter~4.  Since \(U\)
preserves filtered colimits, it follows that
\[
    \Mot_{S,+}(Y)
\]
is compact in \(\HH_\bullet(S)\).

The Thom sphere is the cofiber
\[
    T_S
    \simeq
    \operatorname{cofib}
    \left(
      \Mot_{S,+}(\Gm{}_S)
      \longrightarrow
      \Mot_{S,+}(\mathbb A^1_S)
    \right).
\]
Both \(\Gm{}_S\) and \(\mathbb A^1_S\) are smooth over \(S\) and of finite
presentation.  Their pointed motives are therefore compact.  Compact objects
are closed under finite colimits, so \(T_S\) is compact.

For every smooth \(S\)-scheme \(X\) of finite presentation, distributivity of
the smash product over cofibers gives
\[
\begin{aligned}
    T_S\wedge_S\Mot_{S,+}(X)
    \simeq
    \operatorname{cofib}
    \bigl(
      \Mot_{S,+}(X\times_S\Gm{}_S)
      \longrightarrow
      \Mot_{S,+}(X\times_S\mathbb A^1_S)
    \bigr).
\end{aligned}
\]
The two objects in this cofiber are pointed motives of smooth
finite-presentation \(S\)-schemes and are therefore compact.  Hence
\[
    T_S\wedge_S\Mot_{S,+}(X)
\]
is compact.

Chapter~4 shows that the pointed smooth motives generate
\(\HH_\bullet(S)\) under colimits.  By
Lemma~\ref{lem:compact-closure-generators}, every compact object is obtained
from them by finite colimits and retracts.  The functor
\[
    T_S\wedge_S-:
    \HH_\bullet(S)\longrightarrow\HH_\bullet(S)
\]
preserves finite colimits and retracts and sends each pointed smooth motive to
a compact object.  It therefore preserves every compact object.
\end{proof}

\subsection{Stable generators}

\begin{theorem}[Compact generation of stable motivic spectra]
\label{thm:compact-generation-SH}
The stable presentable category $\SH(S)$ is compactly generated by the objects
\[
    \Sigma^{n\mathcal O_S}
    \Sigma_{T,+}^\infty\Mot_S(X),
    \qquad
    X\in\Sm^{\mathrm{fp}}_{/S},
    \quad n\in\mathbb Z.
\]
Equivalently, writing $p:X\to S$ for the structural morphism, it is generated by
\[
    p_\sharp\Sigma^{n\mathcal O_X}\mathbf 1_X,
    \qquad n\in\mathbb Z.
\]
Every object in the displayed families is compact.
\end{theorem}

\begin{proof}
The category $\HH_\bullet(S)$ is compactly generated by the pointed smooth motives
$\Mot_{S,+}(X)$.  Propositions~\ref{prop:T-symmetric} and \ref{prop:T-compact} verify the hypotheses of
Proposition~\ref{prop:compact-generation-periodization}.  Thus $\SH(S)$ is compactly generated by
positive and negative tensor powers of $T_S$ applied to the suspension spectra of the pointed smooth
motives.  Since tensoring with $T_S$ is $\Sigma^{\mathcal O_S}$, this is the first displayed family.

For $p:X\to S$, Proposition~\ref{prop:p-sharp-on-suspension-spectra} and
Proposition~\ref{prop:p-sharp-virtual-twist} identify
\[
    p_\sharp\Sigma^{n\mathcal O_X}\mathbf 1_X
    \simeq
    \Sigma^{n\mathcal O_S}\Sigma_{T,+}^\infty\Mot_S(X).
\]
This gives the second family.
\end{proof}

\begin{corollary}[Compactness of twisted smooth motives]
\label{cor:twisted-smooth-motives-compact}
Let $p:X\to S$ be smooth of finite presentation and let $\xi\in K(X)$.  Then
\[
    p_\sharp\Sigma^\xi\mathbf 1_X
\]
is compact in $\SH(S)$.
\end{corollary}

\begin{proof}
Tensoring with $\operatorname{th}_X(\xi)$ is an autoequivalence, so it preserves compact objects.
The unit $\mathbf 1_X$ is compact by Theorem~\ref{thm:compact-generation-SH}, applied over $X$.
The right adjoint $p^*$ preserves all colimits, in particular filtered colimits.  Therefore its left
adjoint $p_\sharp$ preserves compact objects.  The result follows.
\end{proof}

\begin{corollary}[Compact sphere]
\label{cor:compact-sphere}
The motivic sphere spectrum $\mathbf 1_S$ is compact.
\end{corollary}

\begin{proof}
Take $X=S$ and $n=0$ in Theorem~\ref{thm:compact-generation-SH}.
\end{proof}

\section{Stable derived nilpotent invariance and classical comparison}
\label{sec:stable-classical-comparison}

\subsection{Compatibility of unstable comparison with Thom spheres}

Let $S_{\mathrm{cl}}$ be the classical truncation of $S$.  Chapter~5 constructed natural symmetric
monoidal equivalences
\[
    \HH_\bullet(S)
    \simeq
    \HH_\bullet(S_{\mathrm{cl}})
\]
for the discrete spectral enhancement and
\[
    \HH_\bullet(S)
    \simeq
    \HH_{\mathrm{cl},\bullet}(S_{\mathrm{cl}})
\]
for the ordinary motivic category.

\begin{proposition}[Compatibility with stabilization objects and Thom twists]
\label{prop:comparison-preserves-T}
Let
\[
    \iota_S:S_{\mathrm{cl}}\longrightarrow S
\]
be the canonical morphism, where \(S_{\mathrm{cl}}\) is regarded as a
discrete spectral scheme.

\begin{enumerate}
    \item Under derived nilpotent invariance,
    \[
        \iota_S^*:
        \HH_\bullet(S)
        \xrightarrow{\simeq}
        \HH_\bullet(S_{\mathrm{cl}}),
    \]
    there are canonical equivalences
    \[
        \iota_S^*T_S
        \simeq
        T_{S_{\mathrm{cl}}},
    \]
    and, for every finite locally free module \(E\) on \(S\),
    \[
        \iota_S^*\Th_S(E)
        \simeq
        \Th_{S_{\mathrm{cl}}}(\iota_S^*E).
    \]

    \item Under ordinary classical comparison,
    \[
        \Theta_S:
        \HH_\bullet(S)
        \xrightarrow{\simeq}
        \HH_{\mathrm{cl},\bullet}(S_{\mathrm{cl}}),
    \]
    there are canonical equivalences
    \[
        \Theta_S(T_S)
        \simeq
        T^{\mathrm{cl}}_{S_{\mathrm{cl}}},
    \]
    and
    \[
        \Theta_S\Th_S(E)
        \simeq
        \Th^{\mathrm{cl}}_{S_{\mathrm{cl}}}
        \bigl(
          \iota_S^*E
        \bigr).
    \]

    \item After periodization, for every
    \[
        \xi\in K(S),
    \]
    there are canonical equivalences
    \[
        \iota_S^*
        \operatorname{th}_S(\xi)
        \simeq
        \operatorname{th}_{S_{\mathrm{cl}}}
        \bigl(
          \iota_S^*\xi
        \bigr)
    \]
    and
    \[
        \Theta_S^{\mathrm{st}}
        \operatorname{th}_S(\xi)
        \simeq
        \operatorname{th}^{\mathrm{cl}}_{S_{\mathrm{cl}}}
        \bigl(
          \iota_S^*\xi
        \bigr).
    \]
    Consequently, for every \(M\in\SH(S)\),
    \[
        \iota_S^*
        \bigl(
          \Sigma^\xi M
        \bigr)
        \simeq
        \Sigma^{\iota_S^*\xi}
        \iota_S^*M
    \]
    and
    \[
        \Theta_S^{\mathrm{st}}
        \bigl(
          \Sigma^\xi M
        \bigr)
        \simeq
        \Sigma^{\iota_S^*\xi}
        \Theta_S^{\mathrm{st}}(M).
    \]
\end{enumerate}
\end{proposition}

\begin{proof}
For derived nilpotent invariance, base change for relative linear spaces gives
\[
    \iota_S^*\mathbb V_S(E)
    \simeq
    \mathbb V_{S_{\mathrm{cl}}}
    \bigl(
      \iota_S^*E
    \bigr).
\]
This equivalence identifies the zero sections and therefore also their open
complements:
\[
\begin{aligned}
    \iota_S^*
    \bigl(
      \mathbb V_S(E)\setminus S
    \bigr)
    \simeq
    \mathbb V_{S_{\mathrm{cl}}}
    \bigl(
      \iota_S^*E
    \bigr)
    \setminus
    S_{\mathrm{cl}}.
\end{aligned}
\]
Since pointed inverse image preserves colimits, it preserves the defining
quotient and gives
\[
    \iota_S^*\Th_S(E)
    \simeq
    \Th_{S_{\mathrm{cl}}}
    \bigl(
      \iota_S^*E
    \bigr).
\]
Taking \(E=\mathcal O_S\) gives
\[
    \iota_S^*T_S
    \simeq
    T_{S_{\mathrm{cl}}}.
\]

For ordinary classical comparison, classical truncation of relative spectra
gives
\[
    \bigl(
      \mathbb V_S(E)
    \bigr)_{\mathrm{cl}}
    \simeq
    \mathbb V_{S_{\mathrm{cl}}}
    \bigl(
      \iota_S^*E
    \bigr).
\]
This equivalence identifies zero sections and their open complements.
The comparison functor of Chapter~5 carries the motive of a smooth spectral
scheme to the ordinary motive of its classical truncation and preserves
pointed colimits.  It therefore carries the quotient defining
\(\Th_S(E)\) to the ordinary Thom quotient:
\[
    \Theta_S\Th_S(E)
    \simeq
    \Th^{\mathrm{cl}}_{S_{\mathrm{cl}}}
    \bigl(
      \iota_S^*E
    \bigr).
\]
Taking \(E=\mathcal O_S\) gives
\[
    \Theta_S(T_S)
    \simeq
    T^{\mathrm{cl}}_{S_{\mathrm{cl}}}.
\]

Both comparison functors are strong symmetric monoidal and preserve the
corresponding stabilization objects.  They therefore periodize to strong
symmetric monoidal stable comparison functors.  For every actual finite
locally free module \(E\), the preceding Thom-space equivalences give
\[
    \iota_S^*
    \operatorname{th}_S(E)
    \simeq
    \operatorname{th}_{S_{\mathrm{cl}}}
    \bigl(
      \iota_S^*E
    \bigr)
\]
and
\[
    \Theta_S^{\mathrm{st}}
    \operatorname{th}_S(E)
    \simeq
    \operatorname{th}^{\mathrm{cl}}_{S_{\mathrm{cl}}}
    \bigl(
      \iota_S^*E
    \bigr).
\]

Pullback of finite locally free modules is exact and defines a natural map
\[
    \iota_S^*:
    K(S)
    \longrightarrow
    K(S_{\mathrm{cl}}).
\]
Naturality of the motivic \(J\)-homomorphism under pullback gives a
coherently commutative square
\[
\begin{tikzcd}[column sep=large,row sep=large]
    K(S) \arrow[r,"J_S"] \arrow[d,"\iota_S^*"']
        & \Pic(\SH(S)) \arrow[d,"\iota_S^*"] \\
    K(S_{\mathrm{cl}})
        \arrow[r,"J_{S_{\mathrm{cl}}}"']
        & \Pic(\SH(S_{\mathrm{cl}})).
\end{tikzcd}
\]
Evaluating at \(\xi\in K(S)\) gives
\[
    \iota_S^*
    \operatorname{th}_S(\xi)
    \simeq
    \operatorname{th}_{S_{\mathrm{cl}}}
    \bigl(
      \iota_S^*\xi
    \bigr).
\]

It remains to prove the corresponding assertion for ordinary classical
comparison.  Let
\[
    J^{\mathrm{cl}}_{S_{\mathrm{cl}}}:
    K(S_{\mathrm{cl}})
    \longrightarrow
    \Pic
    \bigl(
      \SH_{\mathrm{cl}}(S_{\mathrm{cl}})
    \bigr)
\]
be the ordinary motivic \(J\)-homomorphism obtained from the classical
stable Thom construction.

Recall that \(J_S\) is obtained by realizing and looping the simplicial
\(E_\infty\)-map
\[
    \delta_\bullet:
    \bigl(
      S_\bullet\Vect(S)
    \bigr)^\simeq
    \longrightarrow
    B_\bullet\Pic(\SH(S)).
\]
For a filtration
\[
    0=E_0\longrightarrow E_1\longrightarrow\cdots\longrightarrow E_n,
\]
its \(a\)-th entry is the Thom object
\[
    \operatorname{th}_S(E_a/E_{a-1}).
\]

Pullback to \(S_{\mathrm{cl}}\) carries this filtration to
\[
    0=\iota_S^*E_0
    \longrightarrow
    \iota_S^*E_1
    \longrightarrow\cdots\longrightarrow
    \iota_S^*E_n
\]
and identifies its successive quotients:
\[
    \iota_S^*
    \bigl(
      E_a/E_{a-1}
    \bigr)
    \simeq
    \iota_S^*E_a/
    \iota_S^*E_{a-1}.
\]
The actual-bundle Thom comparison therefore gives, levelwise,
\[
\begin{aligned}
    \Theta_S^{\mathrm{st}}
    \operatorname{th}_S
    \bigl(
      E_a/E_{a-1}
    \bigr)
    \simeq
    \operatorname{th}^{\mathrm{cl}}_{S_{\mathrm{cl}}}
    \bigl(
      \iota_S^*E_a/
      \iota_S^*E_{a-1}
    \bigr).
\end{aligned}
\]

The multiplication maps in
\(\delta_\bullet\) are the Thom-additivity equivalences associated with the
exact sequences
\[
    0
    \longrightarrow
    E_b/E_a
    \longrightarrow
    E_c/E_a
    \longrightarrow
    E_c/E_b
    \longrightarrow
    0.
\]
They were constructed from the Cartesian diagrams of relative linear
spaces, zero sections, stable smooth--closed exchange, and the smooth and
closed projection formulas.

Classical comparison carries these relative linear spaces, zero sections,
and Cartesian diagrams to their ordinary classical counterparts.
The unstable comparison transformations of Chapter~5 are compatible with
smooth--closed exchange and the projection transformations, and their
periodizations retain this compatibility.  Hence
\(\Theta_S^{\mathrm{st}}\) carries every spectral Thom-additivity
multiplication map to the corresponding ordinary Thom-additivity
multiplication map.

It follows that the actual-bundle Thom comparisons assemble into a
homotopy-commutative square of simplicial \(E_\infty\)-spaces
\[
\begin{tikzcd}[column sep=large,row sep=large]
    \bigl(
      S_\bullet\Vect(S)
    \bigr)^\simeq
        \arrow[r,"\delta_\bullet"]
        \arrow[d,"\iota_S^*"']
    &
    B_\bullet\Pic(\SH(S))
        \arrow[d,"\Theta_S^{\mathrm{st}}"]
    \\
    \bigl(
      S_\bullet\Vect(S_{\mathrm{cl}})
    \bigr)^\simeq
        \arrow[r,"\delta_\bullet^{\mathrm{cl}}"']
    &
    B_\bullet
    \Pic
    \bigl(
      \SH_{\mathrm{cl}}(S_{\mathrm{cl}})
    \bigr).
\end{tikzcd}
\]
Geometric realization and looping therefore give a coherently commutative
square
\[
\begin{tikzcd}[column sep=large,row sep=large]
    K(S) \arrow[r,"J_S"] \arrow[d,"\iota_S^*"']
        & \Pic(\SH(S))
            \arrow[d,"\Theta_S^{\mathrm{st}}"] \\
    K(S_{\mathrm{cl}})
        \arrow[r,"J^{\mathrm{cl}}_{S_{\mathrm{cl}}}"']
        &
    \Pic
    \bigl(
      \SH_{\mathrm{cl}}(S_{\mathrm{cl}})
    \bigr).
\end{tikzcd}
\]
Evaluating at \(\xi\in K(S)\) gives
\[
    \Theta_S^{\mathrm{st}}
    \operatorname{th}_S(\xi)
    \simeq
    \operatorname{th}^{\mathrm{cl}}_{S_{\mathrm{cl}}}
    \bigl(
      \iota_S^*\xi
    \bigr).
\]

Finally, strong symmetric monoidality gives
\[
\begin{aligned}
    \iota_S^*
    \bigl(
      M\otimes_S\operatorname{th}_S(\xi)
    \bigr)
    &\simeq
    \iota_S^*M
    \otimes_{S_{\mathrm{cl}}}
    \operatorname{th}_{S_{\mathrm{cl}}}
    \bigl(
      \iota_S^*\xi
    \bigr),
\end{aligned}
\]
and
\[
\begin{aligned}
    \Theta_S^{\mathrm{st}}
    \bigl(
      M\otimes_S\operatorname{th}_S(\xi)
    \bigr)
    &\simeq
    \Theta_S^{\mathrm{st}}(M)
    \otimes
    \operatorname{th}^{\mathrm{cl}}_{S_{\mathrm{cl}}}
    \bigl(
      \iota_S^*\xi
    \bigr).
\end{aligned}
\]
These are precisely the two asserted twist equivalences.
\end{proof}

\subsection{Stable nilpotent invariance}

\begin{theorem}[Stable derived nilpotent invariance]
\label{thm:stable-nil-invariance}
The canonical morphism
\[
    \iota_S:
    S_{\mathrm{cl}}
    \longrightarrow
    S,
\]
with \(S_{\mathrm{cl}}\) regarded as a discrete spectral scheme, induces a
symmetric monoidal equivalence
\[
    \iota_S^*:
    \SH(S)
    \xrightarrow{\simeq}
    \SH(S_{\mathrm{cl}}).
\]

For every
\[
    \xi\in K(S)
\]
and every \(M\in\SH(S)\), this equivalence satisfies
\[
    \iota_S^*
    \bigl(
      \Sigma^\xi M
    \bigr)
    \simeq
    \Sigma^{\iota_S^*\xi}
    \iota_S^*M.
\]
No equivalence
\[
    K(S)\simeq K(S_{\mathrm{cl}})
\]
is asserted.
\end{theorem}

\begin{proof}
The pointed unstable inverse-image functor
\[
    \iota_S^*:
    \HH_\bullet(S)
    \longrightarrow
    \HH_\bullet(S_{\mathrm{cl}})
\]
is a symmetric monoidal equivalence by derived nilpotent invariance in
Chapter~5.  Proposition~\ref{prop:comparison-preserves-T} gives
\[
    \iota_S^*T_S
    \simeq
    T_{S_{\mathrm{cl}}}.
\]
Periodizing the unstable equivalence along the corresponding stabilization
objects therefore gives a symmetric monoidal equivalence
\[
    \iota_S^*:
    \SH(S)
    \xrightarrow{\simeq}
    \SH(S_{\mathrm{cl}}).
\]

The formula for virtual Thom twists is
Proposition~\ref{prop:comparison-preserves-T}(3):
\[
    \iota_S^*
    \operatorname{th}_S(\xi)
    \simeq
    \operatorname{th}_{S_{\mathrm{cl}}}
    \bigl(
      \iota_S^*\xi
    \bigr).
\]
Tensoring this equivalence with \(\iota_S^*M\) gives
\[
    \iota_S^*
    \bigl(
      \Sigma^\xi M
    \bigr)
    \simeq
    \Sigma^{\iota_S^*\xi}
    \iota_S^*M.
\]
The construction uses only the pullback map
\[
    K(S)\longrightarrow K(S_{\mathrm{cl}}),
\]
so no \(K\)-theory equivalence is involved.
\end{proof}

\subsection{Ordinary stable motivic comparison}

\begin{definition}[Ordinary stable motivic category]
\label{def:ordinary-SH}
For an ordinary quasi-compact quasi-separated scheme \(S_0\), define the
\(T\)-periodic ordinary motivic category by
\[
    \SH_{\mathrm{cl}}(S_0)
    :=
    \HH_{\mathrm{cl},\bullet}(S_0)
    \bigl[
      (T^{\mathrm{cl}}_{S_0})^{-1}
    \bigr],
\]
where
\[
    T^{\mathrm{cl}}_{S_0}
    :=
    \mathbb A^1_{S_0}/
    \bigl(
      \mathbb A^1_{S_0}\setminus0
    \bigr).
\]
\end{definition}

\begin{proposition}[Comparison with \(\mathbb P^1\)-stabilization]
\label{prop:ordinary-T-P1-comparison}
There is a canonical equivalence in
\(\HH_{\mathrm{cl},\bullet}(S_0)\)
\[
    T^{\mathrm{cl}}_{S_0}
    \simeq
    \bigl(
      \mathbb P^1_{S_0},
      \infty
    \bigr).
\]
Consequently,
\[
    \SH_{\mathrm{cl}}(S_0)
    \simeq
    \HH_{\mathrm{cl},\bullet}(S_0)
    \bigl[
      (\mathbb P^1_{S_0},\infty)^{-1}
    \bigr].
\]
Thus Definition~\ref{def:ordinary-SH} agrees with the conventional
\(\mathbb P^1\)-stable motivic homotopy category.
\end{proposition}

\begin{proof}
Let
\[
    \mathbb A^1_0,
    \qquad
    \mathbb A^1_\infty
\]
be the two standard affine opens of
\(\mathbb P^1_{S_0}\), with
\[
    \mathbb A^1_0\cap\mathbb A^1_\infty
    =
    (\Gm)_{S_0}.
\]
The point at infinity defines a section
\[
    \infty:
    S_0\longrightarrow\mathbb A^1_\infty
    \longrightarrow\mathbb P^1_{S_0}.
\]

In the pointed motivic category there is a natural cofiber sequence
\[
    \bigl(
      \mathbb A^1_\infty,\infty
    \bigr)
    \longrightarrow
    \bigl(
      \mathbb P^1_{S_0},\infty
    \bigr)
    \longrightarrow
    \mathbb P^1_{S_0}/\mathbb A^1_\infty.
\]
Affine-line invariance identifies
\[
    \bigl(
      \mathbb A^1_\infty,\infty
    \bigr)
    \simeq
    0_{S_0}.
\]
Consequently, the quotient morphism induces a canonical equivalence
\[
    \bigl(
      \mathbb P^1_{S_0},\infty
    \bigr)
    \xrightarrow{\simeq}
    \mathbb P^1_{S_0}/\mathbb A^1_\infty.
\]

Nisnevich excision for the standard affine cover identifies the latter
quotient with the corresponding quotient on the other affine chart:
\[
\begin{aligned}
    \mathbb P^1_{S_0}/\mathbb A^1_\infty
    &\simeq
    \mathbb A^1_0/
    \bigl(
      \mathbb A^1_0\cap\mathbb A^1_\infty
    \bigr) \\
    &\simeq
    \mathbb A^1_{S_0}/(\Gm)_{S_0} \\
    &=
    T^{\mathrm{cl}}_{S_0}.
\end{aligned}
\]
Combining the two equivalences gives
\[
    T^{\mathrm{cl}}_{S_0}
    \simeq
    \bigl(
      \mathbb P^1_{S_0},\infty
    \bigr).
\]

Periodization at canonically equivalent objects gives
\[
    \SH_{\mathrm{cl}}(S_0)
    \simeq
    \HH_{\mathrm{cl},\bullet}(S_0)
    \bigl[
      (\mathbb P^1_{S_0},\infty)^{-1}
    \bigr].
\]
\end{proof}


\begin{theorem}[Stable classical comparison]
\label{thm:stable-classical-comparison}
For every quasi-compact quasi-separated spectral scheme \(S\), there is a
natural symmetric monoidal equivalence
\[
    \Theta_S^{\mathrm{st}}:
    \SH(S)
    \xrightarrow{\simeq}
    \SH_{\mathrm{cl}}(S_{\mathrm{cl}}).
\]

The equivalence is compatible with:
\begin{enumerate}
    \item inverse image and tensor products;
    \item smooth direct image;
    \item arbitrary direct image;
    \item stable localization and closed direct image;
    \item exceptional inverse image for closed immersions;
    \item closed base change and smooth--closed exchange;
    \item stable homotopy purity.
\end{enumerate}

For every
\[
    \xi\in K(S)
\]
and every \(M\in\SH(S)\), there is a canonical equivalence
\[
    \Theta_S^{\mathrm{st}}
    \bigl(
      \Sigma^\xi M
    \bigr)
    \simeq
    \Sigma^{\iota_S^*\xi}
    \Theta_S^{\mathrm{st}}(M),
\]
where
\[
    \iota_S^*:
    K(S)
    \longrightarrow
    K(S_{\mathrm{cl}})
\]
is pullback of finite locally free modules.
\end{theorem}

\begin{proof}
The pointed classical-comparison equivalence from Chapter~5 is strong
symmetric monoidal:
\[
    \Theta_S:
    \HH_\bullet(S)
    \xrightarrow{\simeq}
    \HH_{\mathrm{cl},\bullet}(S_{\mathrm{cl}}).
\]
By Proposition~\ref{prop:comparison-preserves-T}, it carries the
stabilization object to the classical stabilization object:
\[
    \Theta_S(T_S)
    \simeq
    T^{\mathrm{cl}}_{S_{\mathrm{cl}}}.
\]
The universal property of periodization therefore gives a natural strong
symmetric monoidal equivalence
\[
    \Theta_S^{\mathrm{st}}:
    \SH(S)
    \xrightarrow{\simeq}
    \SH_{\mathrm{cl}}(S_{\mathrm{cl}}).
\]

We first record the compatibilities obtained by periodization.  The unstable
comparison is compatible with inverse image, pointed tensor products, smooth
direct image, pointed localization, closed direct image, closed base change,
smooth--closed exchange, and homotopy purity.  Each of these functors and
natural transformations is compatible with the corresponding \(T\)-linearity
data.  Proposition~\ref{prop:linear-functor-periodization} and
Proposition~\ref{prop:multilinear-periodization} therefore carry the
unstable comparison transformations to stable comparison transformations.
Since the unstable transformations are equivalences, the stable
transformations are equivalences.

Thus \(\Theta_S^{\mathrm{st}}\) is compatible with:
\[
    f^*,
    \qquad
    p_\sharp,
    \qquad
    \otimes,
    \qquad
    i_*,
\]
as well as stable localization, closed base change, smooth--closed exchange,
and stable homotopy purity.

We next record the compatibilities obtained by adjunction.  Compatibility
with arbitrary direct image follows by taking right adjoints of the
comparison square for inverse image:
\[
    f^*\dashv f_*.
\]
Compatibility with exceptional inverse image for a closed immersion follows
by taking right adjoints of the comparison square for closed direct image:
\[
    i_*\dashv i^!.
\]
Uniqueness of adjoint mates supplies the required coherence with units,
counits, and composition.

It remains to treat Thom twists.  Naturality of the motivic
\(J\)-homomorphism and
Proposition~\ref{prop:comparison-preserves-T} give
\[
    \Theta_S^{\mathrm{st}}
    \operatorname{th}_S(\xi)
    \simeq
    \operatorname{th}^{\mathrm{cl}}_{S_{\mathrm{cl}}}
    \bigl(
      \iota_S^*\xi
    \bigr).
\]
Since \(\Theta_S^{\mathrm{st}}\) is strong symmetric monoidal,
\[
\begin{aligned}
    \Theta_S^{\mathrm{st}}
    \bigl(
      \Sigma^\xi M
    \bigr)
    &=
    \Theta_S^{\mathrm{st}}
    \bigl(
      M\otimes_S\operatorname{th}_S(\xi)
    \bigr) \\
    &\simeq
    \Theta_S^{\mathrm{st}}(M)
    \otimes
    \operatorname{th}^{\mathrm{cl}}_{S_{\mathrm{cl}}}
    \bigl(
      \iota_S^*\xi
    \bigr) \\
    &=
    \Sigma^{\iota_S^*\xi}
    \Theta_S^{\mathrm{st}}(M).
\end{aligned}
\]
This proves the twist formula and completes the comparison.
\end{proof}

\begin{corollary}[Stable comparison of coefficient systems]
\label{cor:stable-coefficient-comparison}
The equivalences $\Theta_S^{\mathrm{st}}$ assemble into an equivalence between the stable brave new
motivic coefficient system and the ordinary stable motivic coefficient system over classical
truncations.
\end{corollary}

\begin{proof}
The coherence follows from the coherent unstable comparison and functoriality of symmetric
monoidal periodization.
\end{proof}

\section{The stable brave new motivic package}
\label{sec:stable-package}

\begin{theorem}[Stable brave new motivic homotopy theory]
\label{thm:stable-brave-new-package}
For every quasi-compact quasi-separated spectral scheme $S$, there is a compactly generated stable
presentably symmetric monoidal $\infty$-category
\[
    \SH(S)
    =
    \HH_\bullet(S)[T_S^{-1}].
\]
The following structures and properties are canonically available.

\begin{enumerate}
    \item There is an adjunction
    \[
        \Sigma_T^\infty:
        \HH_\bullet(S)
        \rightleftarrows
        \SH(S):\Omega_T^\infty,
    \]
    whose left adjoint is colimit-preserving and strong symmetric monoidal.
    Consequently, the right adjoint \(\Omega_T^\infty\) is canonically lax
    symmetric monoidal.

    \item Every morphism $f:T\to S$ induces an adjunction
    \[
        f^*:\SH(S)
        \rightleftarrows
        \SH(T):f_*,
    \]
    where $f^*$ preserves colimits and is strong symmetric monoidal.

    \item Every smooth morphism $p:X\to S$ of finite presentation induces an adjunction
    \[
        p_\sharp:\SH(X)
        \rightleftarrows
        \SH(S):p^*.
    \]

    \item Stable smooth base change holds.  For every Cartesian square
    \[
    \begin{tikzcd}[column sep=large,row sep=large]
        X_T \arrow[r,"g'"] \arrow[d,"p'"']
            & X \arrow[d,"p"] \\
        T \arrow[r,"g"'] & S
    \end{tikzcd}
    \]
    with $p$ smooth of finite presentation, there is a canonical equivalence
    \[
        p'_\sharp(g')^*
        \simeq
        g^*p_\sharp.
    \]

    \item The stable smooth projection formula holds:
    \[
        p_\sharp(E\otimes_Xp^*F)
        \simeq
        p_\sharp(E)\otimes_SF.
    \]

    \item For every complementary closed--open pair
    \[
        Z\xrightarrow{i}S\xleftarrow{j}U,
    \]
    there are natural exact triangles
    \[
        j_\sharp j^*E
        \longrightarrow
        E
        \longrightarrow
        i_*i^*E
    \]
    and
    \[
        i_*i^!E
        \longrightarrow
        E
        \longrightarrow
        j_*j^*E.
    \]

    \item Closed base change, the closed projection formula, and smooth--closed exchange hold:
    \[
        g^*i_*
        \simeq
        i'_*(g')^*,
    \]
    \[
        i_*E\otimes_SF
        \simeq
        i_*(E\otimes_Zi^*F),
    \]
    \[
        p_\sharp k_*
        \simeq
        i_*q_\sharp.
    \]

    \item Stable homotopy purity holds.  For every smooth closed pair
    \[
        Z\lhook\joinrel\longrightarrow X
    \]
    over \(S\), with structural morphism \(z:Z\to S\), there is a canonical
    equivalence
    \[
        \Sigma_{T,S}^\infty
        \bigl(
          X/(X\setminus Z)
        \bigr)
        \simeq
        z_\sharp
        \operatorname{th}_Z
        \bigl(
          N^*_{Z/X}
        \bigr).
    \]

    \item Every finite locally free module \(E\) on \(S\) determines a
    tensor-invertible Thom object
    \[
        \operatorname{th}_S(E)
        \in
        \Pic(\SH(S)).
    \]
    Vector-bundle algebraic \(K\)-theory acts by virtual Thom twists through a
    natural motivic \(J\)-homomorphism
    \[
        J_S:
        K(S)
        \longrightarrow
        \Pic(\SH(S)).
    \]

    \item The category $\SH(S)$ is compactly generated by
    \[
        \Sigma^{n\mathcal O_S}
        \Sigma_{T,+}^\infty\Mot_S(X),
        \qquad
        X\in\Sm^{\mathrm{fp}}_{/S},
        \quad n\in\mathbb Z.
    \]

    \item The theory is derived nilpotent invariant and admits ordinary stable
    comparison:
    \[
        \SH(S)
        \simeq
        \SH(S_{\mathrm{cl}})
        \simeq
        \SH_{\mathrm{cl}}(S_{\mathrm{cl}}).
    \]
    For every \(\xi\in K(S)\), the comparison carries the virtual Thom twist
    \(\Sigma^\xi\) to the twist indexed by its pullback
    \[
        \iota_S^*\xi
        \in
        K(S_{\mathrm{cl}}).
    \]
    No equivalence between \(K(S)\) and \(K(S_{\mathrm{cl}})\) is asserted.
\end{enumerate}

All inverse-image, direct-image, smooth-direct-image, exchange, projection, localization, and
Thom-twist data satisfy coherent composition and identity laws.  Thus
\[
    S\longmapsto\SH(S)
\]
forms a stable $(\ast,\sharp,\otimes)$-formalism.
\end{theorem}

\begin{proof}
Definition and presentable symmetric monoidality are
Definition~\ref{def:SH} and Proposition~\ref{prop:SH-presentable-closed}.  Stability is
Theorem~\ref{thm:SH-stable}.  The suspension adjunction is
Proposition~\ref{prop:SigmaOmegaT-adjunction}.  Arbitrary inverse and direct image are
Proposition~\ref{prop:stable-inverse-image} and Definition~\ref{def:stable-direct-image}, with
coherent functoriality from Corollary~\ref{cor:SH-star-functor} and
Proposition~\ref{prop:stable-direct-image-properties}.

Stable smooth direct image is Theorem~\ref{thm:stable-p-sharp}, with composition from
Proposition~\ref{prop:stable-p-sharp-composition}.  Smooth base change and the smooth projection
formula are Theorems~\ref{thm:stable-smooth-base-change} and
\ref{thm:stable-smooth-projection}.

Stable localization, co-localization, and recollement are
Theorems~\ref{thm:stable-localization}, \ref{thm:stable-colocalization}, and
Proposition~\ref{prop:stable-recollement}.  Closed base change, the closed projection formula, and
smooth--closed exchange are Theorems~\ref{thm:stable-closed-base-change},
\ref{thm:stable-closed-projection}, and \ref{thm:stable-smooth-closed-exchange}.
Stable homotopy purity is
Theorem~\ref{thm:stable-homotopy-purity}.

Thom invertibility is Theorem~\ref{thm:thom-invertibility}; virtual twists are constructed in
Theorem~\ref{thm:motivic-J} and Definition~\ref{def:virtual-thom-suspension}.  Compact generation is
Theorem~\ref{thm:compact-generation-SH}.  Stable nilpotent invariance and ordinary comparison are
Theorems~\ref{thm:stable-nil-invariance} and \ref{thm:stable-classical-comparison}.

Every construction is obtained from coherent unstable geometric data by symmetric monoidal
periodization, adjunction, or passage to the Picard and $K$-theory spaces.  The required coherence
therefore follows from the coherence of those operations.
\end{proof}

The chapter has constructed the stable motivic coefficient system, but no proper or exceptional
functoriality has yet been asserted.  In particular, there is not yet a general exceptional direct image
$f_!$, a general exceptional inverse image $f^!$, proper base change, or smooth proper
ambidexterity.  The next chapter develops spectral projective geometry, stable purity
transformations, smooth proper duality, proper excision, and compactification.  Those results supply
the geometric input from which the exceptional operations will be constructed.
